% Blocks removed from 02-v2 main.tex for the focused version (2026-11-06).
% Full pre-focus manuscript: main_before_focus_59pp.tex in this folder.

%%%%%%%%%%%%%%%%%%%%%%%%%%%%%%%%%%%%%%%%%%%%%%%%%%%%%%%%%%%%%%%%%%%%%%
%% REMOVED: prop:class-undecidable
%%%%%%%%%%%%%%%%%%%%%%%%%%%%%%%%%%%%%%%%%%%%%%%%%%%%%%%%%%%%%%%%%%%%%%
Membership in the class cannot be checked from a grammar; Clark and Eyraud
expected this for substitutability \cite[p.~1743]{ClarkEyraud2007}.

\begin{proposition}[Membership in the class is undecidable]
\label{prop:class-undecidable}
Let $f\ge1$ and let $s$ be a typing on an alphabet $\Sigma$ that contains
distinct letters $\mathtt a,\mathtt b,\mathtt c,\mathtt d$ with
$s(\mathtt a)=s(\mathtt b)$.  It is undecidable, given a context-free grammar
$H$ over $\Sigma$, whether $L(H)\in\Cmcf{f}{s}$; hence also given an $f$-MCFG.
\end{proposition}

\begin{proof}
Put $\Gamma:=\{\mathtt c,\mathtt d\}$.  We reduce universality of context-free
grammars over $\Gamma$, which is undecidable.  Given $H_0$ over $\Gamma$, decide
whether $L(H_0)=\emptyset$; if so, answer ``not universal''.  Otherwise build
a grammar $H$ for $L':=\mathtt a\Gamma^*\cup\mathtt bL(H_0)$.

If $L(H_0)=\Gamma^*$, then $L'=\{\mathtt a,\mathtt b\}\Gamma^*$, a word lies in $L'$
iff its first letter is in $\{\mathtt a,\mathtt b\}$ and no other letter is in
$\Sigma\setminus\Gamma$.  Let $E,F\in\mathcal E_d$ ($d\le f$) and
$\tuple x,\tuple y\in(\Sigma^+)^d$ with $E[\tuple x],E[\tuple y],F[\tuple x]\in L'$.
If $\tuple x$ has a letter outside $\Gamma$, this letter is the only one in
$E[\tuple x]$ and starts it, so $E$ has an empty leftmost fixed word and no letter
outside $\Gamma$; then $E[\tuple y]\in L'$ forces $y_1$ to start with a letter of
$\{\mathtt a,\mathtt b\}$ and $\tuple y$ to have no other letter outside $\Gamma$.
Likewise $F$ has an empty leftmost fixed word and no letter outside $\Gamma$, so
$F[\tuple y]\in L'$.  If $\tuple x\in(\Gamma^+)^d$, the first letter of
$E$ is in $\{\mathtt a,\mathtt b\}$, so $\tuple y\in(\Gamma^+)^d$; the first letter of
$F$ is in $\{\mathtt a,\mathtt b\}$ and $F$ has no other letter outside $\Gamma$, so
$F[\tuple y]\in L'$.  Hence $L'$ is $(f,s)$-tuple-substitutable even for the
trivial typing, and it is regular, so $L'\in\Cmcf{f}{s}$.

Otherwise pick $w\in L(H_0)$ and $w'\in\Gamma^*\setminus L(H_0)$.  The letters
$\mathtt a,\mathtt b$ have the same type and share the context $\blank_1w$, but only
$\mathtt a$ is accepted by $\blank_1w'$; so $L'$ is not
$(1,s)$-tuple-substitutable and $L'\notin\Cmcf{f}{s}$.  Thus
$L(H_0)=\Gamma^*$ iff $L(H)\in\Cmcf{f}{s}$.
\end{proof}

So a learner for $\Cmcf{f}{s}$ cannot test from a grammar whether a target is
in its class; this is the usual situation for distributional learning.


%%%%%%%%%%%%%%%%%%%%%%%%%%%%%%%%%%%%%%%%%%%%%%%%%%%%%%%%%%%%%%%%%%%%%%
%% REMOVED: lem:no-insertion
%%%%%%%%%%%%%%%%%%%%%%%%%%%%%%%%%%%%%%%%%%%%%%%%%%%%%%%%%%%%%%%%%%%%%%
\begin{lemma}[Finite targets admit no typed insertion]
\label{lem:no-insertion}
Let $L$ be $(f,s)$-tuple-substitutable, $1\le d\le f$, $E\in\mathcal E_d$, and
$\tuple x,\tuple y\in(\Sigma^+)^d$ with $s^{(d)}(\tuple x)=s^{(d)}(\tuple y)$ and
$E[\tuple x],E[\tuple y]\in L$.  If $y_h=p_hx_hq_h$ for all $h$, with
$p_1q_1\cdots p_dq_d\ne\eps$, then $E_j[\tuple x]\in L$ for every $j\ge0$, where $E_j$
is $E$ with each $\blank_h$ replaced by $p_h^j\blank_hq_h^j$; these words have
strictly increasing lengths, so $L$ is infinite.
\end{lemma}

\begin{proof}
Each $E_j$ is admissible, since its separators contain those of
$E$.  The tuples share $E$, so $\D_L(\tuple x)=\D_L(\tuple y)$.  If
$E_j[\tuple x]\in L$, then $E_j[\tuple y]\in L$, and $E_j[\tuple y]=E_{j+1}[\tuple x]$.
From $E_0[\tuple x]=E[\tuple x]$, all $E_j[\tuple x]$ lie in $L$, and their lengths
increase strictly.
\end{proof}

In particular, two words of a finite member of $\Cmcf{f}{s}$ that have the same
type are never proper factors of one another ($d=1$, $E=\blank_1$).


%%%%%%%%%%%%%%%%%%%%%%%%%%%%%%%%%%%%%%%%%%%%%%%%%%%%%%%%%%%%%%%%%%%%%%
%% REMOVED: subsec:u1-typings
%%%%%%%%%%%%%%%%%%%%%%%%%%%%%%%%%%%%%%%%%%%%%%%%%%%%%%%%%%%%%%%%%%%%%%
\subsection{Individual Typings onto \texorpdfstring{$U_1$}{U1}}
\label{subsec:u1-typings}

Theorem~\ref{thm:fanout-one-data} concerns pseudovarieties.  For an
individual typing at fan-out one, what matters is which words realize a copy
$\{e,z\}$ of $U_1$: letter-disjoint words of type $e$ force exponential data,
while for the binary typing, whose words of type $1$ are the powers of one
letter, only a quadratic lower bound is known.

\begin{corollary}[A typing-level obstruction at fan-out one]
\label{cor:typing-level-obstruction}
Let $r_{\max}\ge2$ and let $s:\Sigma^+\twoheadrightarrow\mathsf S$ be a typing.
Suppose that distinct idempotents $e,z\in\mathsf S$ satisfy $ez=ze=z$, and that
words $u,v,w\in\Sigma^+$ satisfy $s(u)=s(v)=e$ and $s(w)=z$ and are
\emph{letter-disjoint}: no letter occurs twice in $uvw$.  Let $H_n$ be the image
of the language $L_n$ of Proposition~\ref{prop:three-letter-family}, over three
fresh letters, under $\mathtt a\mapsto u$, $\mathtt b\mapsto v$, $\mathtt c\mapsto w$,
and let $G''_n$ be the image of $G_n$, with rules $S\to E_n$, $E_0\to u$ and
$E_i\to vE_{i-1}E_{i-1}v\mid wE_{i-1}w$.  Then $H_n\in\Cmcf{1}{s}$, $G''_n$ is a
reduced good rank-two presentation of $H_n$ with
\[
 \|G''_n\|=(2|v|+2|w|+9)n+|u|+6,\qquad
 \bar\theta_{G''_n}\le2|u|+2|v|+4n|w|,
\]
and every characteristic sample of $H_n$ for $\mathcal B^{\le r_{\max}}_{1,s}$,
for $\mathcal B_{1,s}$, or for any set-driven operator having characteristic
samples throughout $\Cmcf{1}{s}$ has total encoded length greater than
$(|u|+2|v|)2^n-2|v|$.  So for this single typing no polynomial bounds the
characteristic-data length at fan-out one and rank $r_{\max}$ as in
Definition~\ref{def:uniform-polydata}.
\end{corollary}

\begin{proof}
$L_n$ is finite and Clark--Eyraud substitutable
(Proposition~\ref{prop:three-letter-family}), so $H_n$ is Clark--Eyraud
substitutable by Lemma~\ref{lem:app-u1-code}; it is finite, hence
$H_n\in\Cmcf{1}{s}$.  The grammar is good, and reduced since
$E_i\Rightarrow^*w^iuw^i$.  The size follows from
Definition~\ref{def:standard-mcfg}: the start rule has size $3$, the base rule
$1+|u|$, and the two rules of each level $2|v|+5$ and $2|w|+3$.  Completing every
other choice by the rules $E_j\to wE_{j-1}w$ shows that each rule has a
successful use of length at most $|u|+2n|w|$ (start and unary rules) or
$2|u|+2|v|+(2n+2i-4)|w|$ (the binary rule of level $i$).  In $\{e,z\}$, $e$ is an
identity and $z$ is absorbing, so a word of $H_n$ has type $e$ iff it is the
image of a word of $L_n$ without $\mathtt c$, that is, of $W_n$.  The image $h(W_n)$ has length
$(|u|+2|v|)2^n-2|v|$ by $|h(W_i)|=2|h(W_{i-1})|+2|v|$, and
Proposition~\ref{prop:rare-type} forces it into every such sample.
\end{proof}

With $u,v,w$ letters this is the case of two letters of type $e$; in
particular it applies to every typing onto $U_1$ with at least two letters of
type $1$, since surjectivity gives a letter of type $\zeta$.  It also applies
when no letter has type $e$: for $\mathsf S=\{1,g,z\}$ with identity $1$,
$g^2=1$ and zero $z$, and the typing $\mathtt a,\mathtt b,\mathtt c,\mathtt d\mapsto g$,
$\mathtt t\mapsto z$, take $u=\mathtt{ab}$, $v=\mathtt{cd}$, $w=\mathtt t$ and
$e=1$.  Whether two non-commuting words of type $e$ always suffice, without
letter-disjointness, is open; since two words commute iff they are powers of one word, their
absence means that all words of type $e$ are powers of one primitive word, as
for the binary typing below.

Let $\Sigma=\{\mathtt a,\mathtt b\}$ and let $s_{U,2}:\Sigma^+\twoheadrightarrow U_1$ send $\mathtt a$ to
$1$ and $\mathtt b$ to $\zeta$ (the binary typing).  For this typing, with a single letter of type $1$,
Corollary~\ref{cor:typing-level-obstruction} does not apply; still, its data can
be quadratic in $\bar\theta_G$ for presentations of logarithmic size.

\begin{proposition}[Quadratic data for the binary typing at fan-out one]
\label{prop:binary-u1-quadratic}
For $m,n\ge1$ let $L_{m,n}\subseteq\{\mathtt a,\mathtt b\}^+$ be the set of all
concatenations of $2^n$ words from $\{\mathtt a^{2^m},\mathtt b\}$.
\begin{enumerate}[label=\textup{(\roman*)},leftmargin=*]
\item $L_{m,n}\in\Cmcf{1}{s_{U,2}}$.
\item $L_{m,n}$ has a reduced good fan-out-one presentation $G_{m,n}$ of rule rank
two with $\|G_{m,n}\|=6m+6n+9$ and $\bar\theta_{G_{m,n}}=2^m+2^n-1$.
\item For $\mathcal B_{1,s_{U,2}}$, for $\mathcal B^{\le r}_{1,s_{U,2}}$ \textup($r\ge1$\textup),
and for every set-driven operator with characteristic samples throughout
$\Cmcf{1}{s_{U,2}}$, every characteristic sample $F$ of $L_{m,n}$ has
$\|F\|_+>2^{m+n}$.
\end{enumerate}
Consequently, for $m=n$ one has $\|G_{n,n}\|=12n+9$, $\bar\theta_{G_{n,n}}=2^{n+1}-1$
and $\|F\|_+>\bar\theta_{G_{n,n}}^{\,2}/4$; in particular, for this typing at
fan-out one there is no bound of the form $p(\|G\|)\,\bar\theta_G$ with $p$ a
polynomial, in contrast with Theorem~\ref{thm:local-group-polydata}.
\end{proposition}

The proof is in Appendix~\ref{app:u1-typings}.  The pure word
$\mathtt a^{2^{m+n}}$ is the product of the block length and the number of
blocks, while $\bar\theta_G$ is about their sum, since a shortest sentence using
the long block fills every other position with $\mathtt b$.  Whether the data
for $s_{U,2}$ at fan-out one are polynomial is Open Problem~\ref{op:S1}; at
every fan-out $f\ge2$ they are (Corollary~\ref{cor:two-letter-semilattice}).
Thus, for typings onto $U_1$ at fan-out one, two letters of type $1$, or more
generally two letter-disjoint words of type $1$ next to a letter-disjoint word of
type $\zeta$, force exponential data, and a single letter of type $1$ forces data at least
quadratic in $\bar\theta_G$.


%%%%%%%%%%%%%%%%%%%%%%%%%%%%%%%%%%%%%%%%%%%%%%%%%%%%%%%%%%%%%%%%%%%%%%
%% REMOVED: sec8 unique words, SL reduction, two letters, full content
%%%%%%%%%%%%%%%%%%%%%%%%%%%%%%%%%%%%%%%%%%%%%%%%%%%%%%%%%%%%%%%%%%%%%%
\paragraph{Unique words at fan-out two.}
Theorem~\ref{thm:doubling-no-lift} excludes one family and its extensions.  The next lemma and
proposition constrain every family whose rare-type word is unique, which is the
situation of Proposition~\ref{prop:rare-type} with a single forced word.

\begin{lemma}[Rigidity of a unique word]
\label{lem:unique-word-rigid}
Let $L\in\Cmcf{f}{s}$ and suppose that $W$ is the only word of $L$ of type $m$.
If $1\le d\le f$, $\tuple x,\tuple y\in(\Sigma^+)^d$ have equal types and
$E,F\in\mathcal E_d$ satisfy $E[\tuple x]=E[\tuple y]=F[\tuple x]=W$, then
$F[\tuple y]=W$.  Hence $\{W\}$ is $(f,s)$-tuple-substitutable, and so is $\{v\}$
for every factor $v$ of $W$.
\end{lemma}

\begin{proof}
The tuples share $E$, so $F[\tuple y]\in L$; its type is that of $F[\tuple x]$,
namely $m$, so $F[\tuple y]=W$.  The first claim gives the condition for $\{W\}$.
If $W=pvq$, a violation $(E,\tuple x,\tuple y,F)$ for $\{v\}$ gives the violation
$(pEq,\tuple x,\tuple y,pFq)$ for $\{W\}$; the separators stay nonempty.
\end{proof}

\begin{proposition}[Unique words over two letters of an idempotent type]
\label{prop:unique-two-letter-shapes}
Let $f\ge2$, let $\mathtt a,\mathtt b\in\Sigma$ be distinct letters with
$s(\mathtt a)=s(\mathtt b)=e=e^2$, and let $L\in\Cmcf{f}{s}$.  If $W$ is the only
word of $L$ of type $e$ and $W\in\{\mathtt a,\mathtt b\}^+$ with $|W|=n\ge7$, then,
for some $\{x,y\}=\{\mathtt a,\mathtt b\}$, $W$ is one of
\[
 x^n,\qquad x^{n-1}y,\qquad yx^{n-1},\qquad yx^{n-2}y,\qquad
 xyxy\cdots\ \text{(alternating)}.
\]
Conversely, for each of these words $\{W\}$ is $(f,s)$-tuple-substitutable for
every $f$.  More generally, if $u,v\in\Sigma^+$ are letter-disjoint with
$s(u)=s(v)=e$ and the only word of type $e$ is $W=h(T)$ for the morphism
$h:\mathtt a\mapsto u$, $\mathtt b\mapsto v$ and some $T\in\{\mathtt a,\mathtt b\}^+$ with
$|T|\ge7$, then $T$ has one of these shapes.
\end{proposition}

\begin{proof}
All nonempty words over $\{\mathtt a,\mathtt b\}$ have type $e$, so for
$\{W\}$ the condition of Definition~\ref{def:tuple-substitutable} is that of the
trivial typing.  By Lemma~\ref{lem:unique-word-rigid}, $\{W\}$ is
$2$-substitutable, and Theorem~\ref{thm:app-singleton} in
Appendix~\ref{app:singleton} gives the shapes and the converse.  For the general
statement, a violation $(E,\tuple x,\tuple y,F)$ for $\{T\}$ with the trivial
typing (Appendix~\ref{app:singleton}) is mapped by $h$ to a violation for $\{W\}$:
$h$ keeps the separators nonempty, every nonempty word over $\{u,v\}$ has type
$e$, and $h(F[\tuple y])\ne h(T)$ since $h$ is injective.  This contradicts
Lemma~\ref{lem:unique-word-rigid}, so $T$ passes for arity $2$.
\end{proof}

At fan-out one, the exponential lower bounds come from unique words over two
letters of type $1$ that are far from these shapes
(Proposition~\ref{prop:three-letter-family}, Corollary~\ref{cor:typing-level-obstruction}).
At fan-out at least two, a unique word of such a type is, up to its two end
letters, a power of one letter or alternating.  The proposition does not say
whether such words can be long relative to $\|G\|$ and $\bar\theta_G$: for a
single word, $\bar\theta_G=|W|$, and a target needs other words as well; with
three letters of the type, every factor of $\mathtt{abcabc}\cdots$ also passes
(Lemma~\ref{lem:app-singleton-periodic}),
and for a non-idempotent type the classification changes (for
$\mathtt a,\mathtt b\mapsto g$ with $g^2\ne g$, $\{\mathtt{aaabbaba}\}$ passes for
every $f$).

\paragraph{Reduction to semilattices.}
Let $\mathbf{SL}$ be the pseudovariety of finite semilattices (commutative
idempotent semigroups).  For a nonempty finite alphabet $\Sigma$, the
\emph{content typing} $c_\Sigma$ maps $w\in\Sigma^+$ to the set of letters
occurring in $w$, a morphism onto the nonempty subsets of $\Sigma$ under union.

\begin{lemma}[Coarser typings inherit characteristic samples]
\label{lem:typing-monotone}
Let $s\preceq t$ be typings, $f,r\ge1$, and $L\in\Cmcf{f}{s}$.  For every
finite $K\subseteq L$,
\[
 L\bigl(\mathcal B^{\le r}_{f,t}(K)\bigr)\subseteq
 L\bigl(\mathcal B^{\le r}_{f,s}(K)\bigr)\subseteq L,
\]
and likewise for $\mathcal B_{f,t}$ and $\mathcal B_{f,s}$.  Hence every
characteristic sample of $L$ for $\mathcal B^{\le r}_{f,t}$ (resp.\
$\mathcal B_{f,t}$) is one for $\mathcal B^{\le r}_{f,s}$ (resp.\
$\mathcal B_{f,s}$).
\end{lemma}

\begin{proof}
Both grammars have the same nonterminals and the same Start, Const and Comp
rules (Definition~\ref{def:reconstruction}).  A Link rule of the $t$-grammar
joins observed tuples with a common context and equal $t$-types; since
$\ker t\subseteq\ker s$, their $s$-types are equal, so it is a Link rule of the
$s$-grammar.  This gives the first inclusion; the second is
Proposition~\ref{prop:reconstruction-soundness}.  If $F$ is characteristic for
the $t$-operator and $F\subseteq K\subseteq L$, the left-hand side equals $L$,
hence so does the middle term.
\end{proof}

\begin{proposition}[Reduction of the fan-out boundary to semilattices]
\label{prop:sl-reduction}
Fix $f\ge1$ and $r_{\max}\ge1$, and read uniform polynomial characteristic-data
length at fan-out $f$ and rank $r_{\max}$
(Definition~\ref{def:uniform-polydata}).
\begin{enumerate}[label=\textup{(\roman*)},leftmargin=*]
\item Every pseudovariety $\mathbf V\not\subseteq\mathbf{LG}$ contains
$\mathbf{SL}$.  Consequently, the pseudovarieties with uniform polynomial
characteristic-data length are exactly those contained in $\mathbf{LG}$ if and
only if $\mathbf{SL}$ does not have it.
\item $\mathbf{SL}$ has uniform polynomial characteristic-data length if and
only if, for every nonempty finite alphabet $\Sigma$, the content typing
$c_\Sigma$ satisfies the condition of Definition~\ref{def:uniform-polydata}.
\end{enumerate}
\end{proposition}

\begin{proof}
(i) A finite semigroup $\mathsf S\notin\mathbf{LG}$ is the positive image of
its evaluation typing over the alphabet $\mathsf S$, so it has a subsemigroup
isomorphic to $U_1$ (Theorem~\ref{thm:lg-characterization}); hence
$U_1\in\mathbf V$.  Every finite semilattice $\mathsf H$ embeds into
$U_1^{\mathsf H}$ via $x\mapsto(\chi_z(x))_{z\in\mathsf H}$, where $\chi_z(x)=1$ if
$zx=z$ and $\chi_z(x)=\zeta$ otherwise (a morphism, since $zxy=z$ iff $zx=z$ and
$zy=z$), so $\mathbf{SL}\subseteq\mathbf V$.  The property of
Definition~\ref{def:uniform-polydata} passes from a pseudovariety to its
subpseudovarieties, because fewer typings are quantified.  If $\mathbf{SL}$
lacks it, so does every $\mathbf V\not\subseteq\mathbf{LG}$, while every
$\mathbf V\subseteq\mathbf{LG}$ has it by Theorem~\ref{thm:local-group-polydata}.
Conversely, $U_1\notin\mathbf{LG}$, so $\mathbf{SL}\not\subseteq\mathbf{LG}$.

(ii) The image of $c_\Sigma$ is a semilattice, which gives one direction.  For
the other, let $s:\Sigma^+\twoheadrightarrow\mathsf S$ with $\mathsf S\in\mathbf{SL}$.
Since $\mathsf S$ is commutative and idempotent, $s(w)$ is the product of
$s(a)$ over the letters $a$ occurring in $w$, so $s\preceq c_\Sigma$ and
$\Cmcf{f}{s}\subseteq\Cmcf{f}{c_\Sigma}$ (item~(a) after
Definition~\ref{def:class-cmcf}).  For $L\in\Cmcf{f}{s}$ and a presentation $G$,
a characteristic sample for $\mathcal B^{\le r_{\max}}_{f,c_\Sigma}$ of length at
most $p_{c_\Sigma}(\|G\|,\bar\theta_G)$ is one for
$\mathcal B^{\le r_{\max}}_{f,s}$ by Lemma~\ref{lem:typing-monotone}; take
$p_{s}:=p_{c_\Sigma}$.
\end{proof}

By Theorem~\ref{thm:fanout-one-data}, $\mathbf{SL}$ lacks the property at
fan-out one when $r_{\max}\ge2$; at fan-out two the question is open (Open
Problem~\ref{op:fanout-two}), and the next theorem settles two-letter
alphabets at every fan-out $f\ge2$.  Rare-type lower bounds transfer to content
typings: for $s\preceq c_\Sigma$ every content type contained in an $s$-type $m$
has shortest word at least $\mu_m(L,s)$.

\paragraph{Two letters.}
Call $L$ \emph{Parikh-saturated} if membership of a word of $\Sigma^+$ in $L$
depends only on its Parikh vector.

\begin{theorem}[Content typings on two letters at every fan-out $f\ge2$]
\label{thm:two-letter-content}
Let $\Sigma=\{\mathtt a,\mathtt b\}$, $c:=c_\Sigma$, $f\ge2$, $r_{\max}\ge2$, and put
$D_f:=(6f+1)3^f$ and
$C_{f,r_{\max}}:=(r_{\max}+1)\bigl(4r_{\max}(3f+1)+D_f\bigr)+3$, and let
$H_m:=\sum_{k=1}^m1/k\le1+\ln m$.  For every
$L\in\Cmcf{f}{c}$ and every reduced good fan-out-$\le f$, rank-$\le r_{\max}$
presentation $G$ of $L\cap\Sigma^+$,
$\tau(G,c)\le C_{f,r_{\max}}|V|^2\bar\theta_G^{\,4}H_{\bar\theta_G}$, and the operators
$\mathcal B^{\le r_{\max}}_{f,c}$ and $\mathcal B_{f,c}$ have a common
characteristic sample $K$ with
\[
 \|K\|_+\le4^{fr_{\max}}\bigl(3+r_{\max}C_{f,r_{\max}}\bigr)\|G\|^3\bar\theta_G^{\,4}H_{\bar\theta_G}
 +\mathbf 1_{\{\eps\in L\}}.
\]
If $L$ is not Parikh-saturated, then
$t_{\mathrm{typ}}(G,c)\le D_f(r_{\max}+1)|V|\bar\theta_G$ and one may take
$\|K\|_+\le4^{fr_{\max}}\bigl(3+r_{\max}D_f(r_{\max}+1)\bigr)\|G\|^2\bar\theta_G
+\mathbf 1_{\{\eps\in L\}}$.  For $f=2$, $D_2=117$ and
$C_{2,r_{\max}}=28r_{\max}(r_{\max}+1)+117(r_{\max}+1)+3$.
\end{theorem}

\begin{corollary}[Semilattice typings on two letters]
\label{cor:two-letter-semilattice}
Let $f\ge2$, $r_{\max}\ge2$ and let $s:\{\mathtt a,\mathtt b\}^+\twoheadrightarrow\mathsf S$ be a
typing with $\mathsf S\in\mathbf{SL}$, for instance $s_{U,2}$
\textup($\mathtt a\mapsto1$, $\mathtt b\mapsto\zeta$\textup).  Then $s$ satisfies
the condition of Definition~\ref{def:uniform-polydata} at fan-out $f$ and rank
$r_{\max}$, with the bound of Theorem~\ref{thm:two-letter-content}, for
$\mathcal B^{\le r_{\max}}_{f,s}$ and for $\mathcal B_{f,s}$.
\end{corollary}

\begin{proof}
As in the proof of Proposition~\ref{prop:sl-reduction}(ii), $s\preceq c$ and
$\Cmcf{f}{s}\subseteq\Cmcf{f}{c}$; apply Theorem~\ref{thm:two-letter-content} and
Lemma~\ref{lem:typing-monotone}.
\end{proof}

The proof is in Appendix~\ref{app:two-letter-content}.  It uses the
conditions of arities one and two, which hold at every fan-out $f\ge2$, through
runs of one letter: once $x^4$ occurs in a word of
$L$ for a letter $x$, membership depends only on the $x$-skeleton (each maximal
run of $x$ replaced by one $x$) and the number of letters $x$, and two runs of the other letter next to a run $xx$ either
force Parikh saturation or merge into a forbidden longer run.  Without
saturation this bounds the antichains of heavy nodes in a shortest typed
derivation; with saturation, node-minimal derivations of words with a
prescribed Parikh vector are analysed directly.  The theorem does not separate
fan-outs: whether the two-letter content typing needs superpolynomial data at
fan-out one is open, and by Lemma~\ref{lem:typing-monotone} a negative answer
to Open Problem~\ref{op:S1} would imply that $c_{\{\mathtt a,\mathtt b\}}$ needs
superpolynomial data at fan-out one, since $s_{U,2}\preceq c_{\{\mathtt a,\mathtt b\}}$.  For three or more
letters a long component of content $\{\mathtt a,\mathtt b\}$ need not contain a
long run, and the question is open.

\paragraph{Components containing every letter.}
One part of the question is settled for every alphabet and every fan-out, by a
property of grammars alone.

\begin{proposition}[Full-content types]
\label{prop:full-content-types}
Let $\Sigma$ be a nonempty finite alphabet, $m:=|\Sigma|$, $c:=c_\Sigma$, and let
$G=(V,\Sigma,P,S,\fan)$ be reduced and good, of fan-out $\le f$ and rank $\le r$.
Every typed nonterminal $X$ of arity $d$ for $c$ whose components all have type
$\Sigma$ satisfies
\[
 \min_{\tuple u\in L_X}\|\tuple u\|_1\le md(2^m-1)^f(r+1)|V|\bar\theta_G .
\]
If $L(G)\in\Cmcf{f}{c}$, then $\tau_X$ satisfies the same bound.  If some word of
$L(G)$ contains every letter of $\Sigma$, then
$\mu_\Sigma(L(G))\le m(2^m-1)^f(r+1)|V|\bar\theta_G$.
\end{proposition}

\begin{proof}
Put $\theta:=\bar\theta_G$.  Fix $\tuple u\in L_X$ of least total length and a
derivation $T$ of $\tuple u$ from $X$ in $\widetilde G_0$ with the fewest nodes.
In every component of $\tuple u$ mark one occurrence of each letter; there are at
most $md$ marks.  Call a node of $T$ heavy if its yield has total length greater
than $\theta$.  Suppose that the yield of a heavy node $u$, with $G$-nonterminal
$B$, contains no mark.  By Lemma~\ref{lem:thickness-equivalence}(ii) some
$\omega\in L_B(G)$ has $\|\omega\|_1\le\theta$; replacing the subtree at $u$, with
the indices erased, by a $G$-derivation of $\omega$ gives a $G$-derivation from
$A$ of a tuple $\tuple u'$ that is shorter than $\tuple u$ and still contains all
marked letters, each in its component.  So every component of $\tuple u'$ has
type $\Sigma$, and by Proposition~\ref{prop:output-invariants}(ii) the new
derivation lifts to one from $X$.  Inserting it at an occurrence of $X$ in a
successful derivation of $\widetilde G_0$ gives a successful derivation, so its
nonterminals and rules belong to $\widetilde G_0$ and $\tuple u'\in L_X$,
contradicting the choice of $\tuple u$.  Hence the yield of every heavy node
contains a mark.  Heavy nodes at one depth have disjoint yields, so there are at
most $md$ of them.  No typed nonterminal repeats on a path of $T$, since
replacing the subtree at the upper occurrence by the one at the lower gives a
derivation from $X$ whose yield is not longer and which has fewer nodes; there
are at most $(2^m-1)^f|V|$ typed nonterminals, so $T$ has at most
$md(2^m-1)^f|V|$ heavy nodes.  If the root is not heavy,
$\|\tuple u\|_1\le\theta$.  Otherwise every letter of $\tuple u$ is written by a
heavy node or lies in the yield of a non-heavy child of one; each heavy node
writes at most $\theta$ letters (Lemma~\ref{lem:thickness-equivalence}(i)) and has
at most $r$ children, which gives the bound.  The bound on $\tau_X$ follows from
$L_X\subseteq\operatorname{Cl}_{L,c}(X)$ (Lemma~\ref{lem:typed-homogeneity}), and
the bound on $\mu_\Sigma$ is the case $X=S_{(\Sigma)}$, $d=1$.
\end{proof}

So for three letters at fan-out two, typed nonterminals all of whose components
contain $\mathtt a$, $\mathtt b$ and $\mathtt c$ have representatives of length at
most $294(r+1)|V|\bar\theta_G$.  By
Corollary~\ref{cor:distribution-representatives}, superpolynomial data for
$c_\Sigma$ need a typed nonterminal $X$ with superpolynomial $\tau_X$ and a
component that misses a letter.


%%%%%%%%%%%%%%%%%%%%%%%%%%%%%%%%%%%%%%%%%%%%%%%%%%%%%%%%%%%%%%%%%%%%%%
%% REMOVED: old conclusion
%%%%%%%%%%%%%%%%%%%%%%%%%%%%%%%%%%%%%%%%%%%%%%%%%%%%%%%%%%%%%%%%%%%%%%
\section{Conclusion and Open Problems}
\label{sec:conclusion}

For a reduced good presentation $G$ and fixed alphabet, fan-out, rank and typing,
Theorem~\ref{thm:data-from-typed-thickness} reduces the total encoded length of
characteristic samples to the shortest typed yield $t_{\mathrm{typ}}(G,s)$, and
Corollary~\ref{cor:distribution-representatives} further to the representative
length $\tau(G,s)$.  Owner-wise simultaneous replacement gives
$O(\|G\|^2\bar\theta_G)$ for window--group and local-group typings at every
fan-out (Theorem~\ref{thm:A}(i)), and at every fan-out the product
$\|G\|\bar\theta_G$ is necessary under every typing over a sufficiently large
alphabet (Theorem~\ref{thm:A}(ii)).  At
fan-out one, $\mathbf{LG}$ is the exact boundary for pseudovarieties, also for
identification in polynomial time and thick data, and outside it the failure is
exponential (Theorem~\ref{thm:B}, Corollary~\ref{cor:iptd}).  At fan-out two the
exponential fan-out-one family does not lift (Theorem~\ref{thm:C}).

\paragraph{Typings onto $U_1$.}
Every finite semigroup outside $\mathbf{LG}$ contains a copy of $U_1$
(Theorem~\ref{thm:lg-characterization}), so the cost of individual typings
outside $\mathbf{LG}$ is decided first for typings onto $U_1$.  At fan-out one,
two letters of type $1$ force exponential data, and so do letter-disjoint words
of the two types in any typing
(Corollary~\ref{cor:typing-level-obstruction}); and the binary typing $s_{U,2}$,
with a single letter of type $1$, needs data quadratic in $\bar\theta_G$ for
presentations of logarithmic size (Proposition~\ref{prop:binary-u1-quadratic}).
At every fan-out $f\ge2$, $s_{U,2}$ admits data polynomial in $\|G\|$ and
$\bar\theta_G$ (Corollary~\ref{cor:two-letter-semilattice}).

\begin{openproblem}[The binary typing at fan-out one]
\label{op:S1}
For $s_{U,2}$, fan-out one and $r_{\max}\ge2$: is the characteristic-data length
bounded by a polynomial in $\|G\|$ and $\bar\theta_G$?  By
Proposition~\ref{prop:binary-u1-quadratic} such a polynomial has degree at least
two in $\bar\theta_G$.  By Lemma~\ref{lem:no-insertion}, a finite target has
at most one word of type $1$, since such words are powers of $\mathtt a$.  Together with Corollary~\ref{cor:two-letter-semilattice}, a
negative answer would show that one typing needs superpolynomial data at
fan-out one and polynomial data at fan-out two.
\end{openproblem}

\begin{openproblem}[Semilattices at fan-out two]
\label{op:fanout-two}
Fix $r_{\max}\ge2$.  Does $\mathbf{SL}$ have uniform polynomial
characteristic-data length at fan-out two, that is
(Proposition~\ref{prop:sl-reduction}(ii)), does every content typing
$c_\Sigma$ admit, for every $L\in\Cmcf{2}{c_\Sigma}$ and every reduced good
fan-out-$\le2$, rank-$\le r_{\max}$ presentation $G$ of $L\cap\Sigma^+$,
characteristic samples for $\mathcal B^{\le r_{\max}}_{2,c_\Sigma}$ of length
polynomial in $\|G\|$ and $\bar\theta_G$?  By
Proposition~\ref{prop:sl-reduction}(i), a negative answer makes $\mathbf{LG}$
the pseudovariety boundary at fan-out two as well, and a positive answer
shows that the boundary depends on the fan-out.  For two-letter alphabets the
answer is positive, at every fan-out $f\ge2$
(Theorem~\ref{thm:two-letter-content}).  For a negative answer, by
Proposition~\ref{prop:rare-type} a family in $\Cmcf{2}{c_\Sigma}$ with
$\mu^\star(L_n,c_\Sigma)=n^{\omega(1)}$ and $\|G_n\|\bar\theta_{G_n}=n^{O(1)}$
would suffice; Theorem~\ref{thm:doubling-no-lift} excludes the doubling family.  By Proposition~\ref{prop:full-content-types}, the type $\Sigma$ cannot be the long type of such a family.  For a single typing onto $U_1$, polynomial data at fan-out
two with two letters of type $1$ would show that one typing needs exponential
data at fan-out one and polynomial data at fan-out two.
\end{openproblem}

Both open problems meet in one question about words that are the only ones of
their type.  Fix $r_{\max}\ge2$.  Say that $(f,s)$ has \emph{short unique words}
if there is a polynomial $p$ such that $|W|\le p(\|G\|,\bar\theta_G)$ whenever
$L\in\Cmcf{f}{s}$, $G$ is a reduced good fan-out-$\le f$, rank-$\le r_{\max}$
presentation of $L\cap\Sigma^+$, and $W$ is the only word of $L$ of its type.

\begin{proposition}[The rare-type route]
\label{prop:unique-word-route}
\begin{enumerate}[label=\textup{(\roman*)},leftmargin=*]
\item If $(1,s_{U,2})$ does not have short unique words, Open
Problem~\ref{op:S1} has a negative answer, and the long unique words are
powers of $\mathtt a$.
\item If $(2,s_U)$, for the typing $s_U$ of
Proposition~\ref{prop:three-letter-family}, does not have short unique words,
Open Problem~\ref{op:fanout-two} has a negative answer, and $\mathbf{LG}$ is the
pseudovariety boundary at fan-out two; the long unique words of type $1$ have the
shapes of Proposition~\ref{prop:unique-two-letter-shapes}.
\item $(f,s_{U,2})$ has short unique words for every $f\ge2$.
\end{enumerate}
\end{proposition}

\begin{proof}
By Proposition~\ref{prop:rare-type}, every characteristic sample of $L$ contains
$W$, so long unique words give long characteristic data, and short
characteristic data give short unique words.  (i) The words of $s_{U,2}$-type
$1$ are the powers of $\mathtt a$.  (ii) The image $U_1$ of $s_U$ is a semilattice,
so $\mathbf{SL}$ fails to have uniform polynomial data at fan-out two, and
Proposition~\ref{prop:sl-reduction}(i) applies; the words of $s_U$-type $1$ lie in
$\{\mathtt a,\mathtt b\}^+$, and $s_U(\mathtt a)=s_U(\mathtt b)=1$ is idempotent.  (iii)
Corollary~\ref{cor:two-letter-semilattice}.
\end{proof}

So a negative answer to either problem along the rare-type route needs a family
of targets whose unique word of some type is, essentially, a power, as at fan-out
one in Proposition~\ref{prop:binary-u1-quadratic}, but superpolynomially long; by
(iii) such a family must use a third letter at fan-out two.  Unique words are
not the only route: a target may force long data through many words of one type.

Such a family also needs long words of the other type, whatever its fan-out
and even without substitutability.  The reason is a property of grammars alone.

\begin{lemma}[Long unique words need long companions]
\label{lem:unique-word-gap}
Let $G=(V,\Sigma,P,S,\fan)$ be reduced and good of rank at most $r$, and let
$W\in L(G)$.
\begin{enumerate}[label=\textup{(\roman*)},leftmargin=*]
\item If every other word of $L(G)$ has length at most $M$, then
$|W|\le\bar\theta_G+(r+1)M$.
\item If $W$ is the only word of $L(G)$ in $\Gamma^+$ for some $\Gamma\subseteq\Sigma$,
and $|W|>\bar\theta_G$, then some word of $L(G)$ contains at least
$|W|/((r+1)|V|\bar\theta_G)$ letters outside $\Gamma$.
\end{enumerate}
For a typing onto $U_1$ and the only word $W$ of type $1$, (i) applies with $M$ the
greatest length of a word of type $\zeta$, and (ii) with $\Gamma$ the set of letters
of type $1$.
\end{lemma}

\begin{proof}
(i) 
Put $N:=|W|$ and $\theta:=\bar\theta_G$, assume $N>\theta+(r+1)M$, and fix a
derivation $T$ of $W$; let $t_v$ be the total length of the tuple derived at a
node $v$, with nonterminal $A_v$.  Suppose that $M<t_v<N-M$.  Replacing the
subtree at $v$ by a derivation of any $\tuple z\in L_{A_v}(G)$ gives a word of
length $N-t_v+\|\tuple z\|_1>M$ by nondeletion, so it is $W$, and
$\|\tuple z\|_1=t_v$.  Then every successful derivation using the rule
at $v$ derives there a tuple of length $t_v>M$, so its word is $W$; the rule thickness of that rule is $N>\theta$, a contradiction.  So no node
has $M<t_v<N-M$.  The root has $t=N\ge N-M$.  A node with $t_v\ge N-M>\theta+rM$
writes at most $\theta$ letters (Lemma~\ref{lem:thickness-equivalence}(i)) and has
at most $r$ children, so some child has $t_u>M$, hence $t_u\ge N-M$.  This gives
an infinite path in the finite tree $T$.

(ii) Let $T$ be a derivation of $W$ with the fewest nodes.  No nonterminal
repeats on a path of $T$: replacing the subtree at the upper occurrence by the one
at the lower occurrence gives a word of $L(G)\cap\Gamma^+$, hence $W$, with fewer
nodes.  So
$T$ has height at most $|V|$.  Call a node heavy if its tuple has total length
greater than $\theta:=\bar\theta_G$.  Heavy nodes are closed under ancestors, the
root is heavy, and every letter of $W$ is written by a heavy node or lies in a
non-heavy child of one, so there are at least $|W|/((r+1)\theta)$ heavy nodes,
and at least $|W|/((r+1)|V|\theta)$ of them at one depth.  For each such node $u$
fix a tuple $\omega_u\in L_{A_u}(G)$ with $\|\omega_u\|_1\le\theta$
(Lemma~\ref{lem:thickness-equivalence}(ii)); it contains a letter outside
$\Gamma$, since otherwise replacing the subtree at $u$ would give a shorter word of
$L(G)\cap\Gamma^+$.
Replacing all these pairwise incomparable subtrees at once gives a word of $L(G)$
that keeps all these letters, by nondeletion.
\end{proof}

In particular, a family as in Proposition~\ref{prop:unique-word-route}(i)
or~(ii) needs words of type $\zeta$ that are superpolynomially long and, by (ii),
one word with superpolynomially many letters of type $\zeta$; it cannot be built
from a long unique word and words with few letters of the other type.

The sub-targets used so far delete a type.  Proposition~\ref{prop:abelian-coset-filter}
offers others, $L\cap\varphi^{-1}(\gamma+H)$ for a morphism $\varphi$ into a finite
abelian group, whose complements in $L$ may contain many words of many types.
Through Lemma~\ref{lem:distinguishing} they force only short words, at every
fan-out and under every typing.

\begin{proposition}[Abelian filters are detected by short words]
\label{prop:abelian-filter-short}
Let $G=(V,\Sigma,P,S,\fan)$ be reduced and good and $L=L(G)$.  There is a set
$K\subseteq L$ of at most $2|P|+1$ words, each of length at most $\bar\theta_G$, with
the following property: for every monoid morphism $\varphi:\Sigma^*\to\mathsf A$
into an abelian group, $H\le\mathsf A$ and $\gamma\in\mathsf A$, if some word of $L$
lies outside $\varphi^{-1}(\gamma+H)$, then some word of $K$ does.  In particular the
differences of the Parikh vectors of the words of $K$ generate the group generated
by those of $L$.
\end{proposition}

\begin{proof}
Write $\varphi$ additively; it factors through the Parikh map $\pi$.  For each
nonterminal $B$ fix a yield $u_B$ of least total length, so $\|u_B\|_1\le\bar\theta_G$
(Lemma~\ref{lem:thickness-equivalence}(ii)), and let $q_B$ be its total Parikh
vector.  For a rule $R:A\to\rho(B_1,\ldots,B_k)$ writing terminals of Parikh vector
$t_R$ put $\delta_R:=t_R+\sum_jq_{B_j}-q_A$.  In any derivation of a word $w$ the
vectors $q$ of non-root nodes cancel, so
\[
 \pi(w)=q_S+\sum_v\delta_{R(v)},
\]
the sum over the nodes $v$ with rules $R(v)$.  Take a successful derivation of
length $\theta_G(R)$ using $R$, and replace at one use of $R$ either the whole
$A$-subtree by a derivation of $u_A$, or only the child subtrees by derivations
of the $u_{B_j}$; by nondeletion and minimality the two words $p,p'\in L$ obtained
have length at most $\theta_G(R)$, and $\pi(p')-\pi(p)=\delta_R$; put them into $K$,
together with the word $u_S$.  Now suppose that all words of $K$ lie in
$\varphi^{-1}(\gamma+H)$.  Then $\varphi(u_S)\in\gamma+H$ and
$\varphi(\delta_R)=\varphi(p')-\varphi(p)\in H$ for every $R$, so
$\varphi(w)\in\gamma+H$ for every $w\in L$.  For the last claim take
$\varphi=\pi$, $\gamma=\pi(u_S)$ and $H$ the group generated by the differences
of the Parikh vectors of the words of $K$.
\end{proof}

The bound is attained by $\{\mathtt a^m,\mathtt b^m\}$ with the filter
$|w|_{\mathtt b}\equiv0\pmod2$.  Intersecting with a set of types gives nothing
beyond deleting one type: $L\setminus\widehat s^{-1}(\mathsf{Acc})$ is the union of
the sets of words of $L$ of the types outside $\mathsf{Acc}$, so its shortest word
has length $\mu_m(L)$ for one of these types $m$, and deleting $m$ alone already
forces a word of length at least $\mu_m(L)$.  So along Lemma~\ref{lem:distinguishing}, a negative
answer to Open Problem~\ref{op:S1} or~\ref{op:fanout-two} needs sub-targets that are
not abelian filters: deleting a type, which is constrained by
Lemma~\ref{lem:unique-word-rigid}, Proposition~\ref{prop:unique-two-letter-shapes}
and Lemma~\ref{lem:unique-word-gap}, or sub-targets that no abelian invariant
separates.

For a different quantity the arity does matter: for the minimum size of a
finite-monoid observation that preserves guarded tuple substitution through
arity $f$, the least arity at which this minimum stabilizes uniformly over a
pseudovariety of finite monoids is $1$, $2$ or $\infty$
\cite{KuriyamaSCL2026}.  Whether the characteristic-data boundary depends on
the fan-out is the corresponding question for data length.

\paragraph{Tightness of the exponents.}
In $O(\|G\|^2\bar\theta_G)$, for window--group and local-group typings alike,
one factor $\|G\|$ counts the witness words and the other bounds root-to-leaf
paths in a $\lambda$-minimal derivation by the number of node labels
(Lemma~\ref{lem:heavy-node-count}).  The path factor is genuine for the witness
construction: the grammars $S\to SS\mid\mathtt a\mid\mathtt b\mid A_0\mid\cdots\mid A_n$,
$A_0\to\mathtt b$, $A_i\to\mathtt a\mid A_{i-1}\mathtt b$ ($1\le i\le n$) generate
$\{\mathtt a,\mathtt b\}^+$ with $\|G_n\|=10n+16$ and $\bar\theta_{G_n}=2$, but
$\mathtt b^{n+1}$ is the only yield of $A_n$ of $s_{1,1}$-type
$\langle\mathtt b,\mathtt b\rangle$, so $t_{\mathrm{typ}}(G_n,s_{1,1})\ge n+1$.  For this
target, however, $\tau(G_n,s_{1,1})\le2$
(Corollary~\ref{cor:distribution-representatives}), so the example bounds the
witness method, not the data.  For the trivial typing the upper bound
$O(\|G\|\bar\theta_G)$ is attained by the marker families
(Corollary~\ref{cor:trivial-exact}); for
every typing with $\mathsf S\in\mathbf{LG}$ the remaining gap is a single factor
$\|G\|$, a gap between the known examples and the bounds rather than a proved
ratio.  Neither lower-bound argument closes it: for a fixed typing the
rare-type argument yields at most $|\mathsf S|(t_{\mathrm{typ}}+1)$, and the marker
families consist of $O(\|G\|)$ words of length $\bar\theta_G=t_{\mathrm{typ}}$.
For noncommutative
groups the constant $M_{\mathbb G,f}$ is doubly exponential in $f$ because of the
Ramsey-type homogenization (Lemma~\ref{lem:order-homogenize}); whether
$\mathrm{poly}(|\mathbb G|^f)$ owners suffice, as for abelian groups, is open.


%%%%%%%%%%%%%%%%%%%%%%%%%%%%%%%%%%%%%%%%%%%%%%%%%%%%%%%%%%%%%%%%%%%%%%
%% REMOVED: app:u1-typings
%%%%%%%%%%%%%%%%%%%%%%%%%%%%%%%%%%%%%%%%%%%%%%%%%%%%%%%%%%%%%%%%%%%%%%
\section{Typings onto \texorpdfstring{$U_1$}{U1} at Fan-Out One}
\label{app:u1-typings}

\subsection{Letter-Disjoint Images of Finite Substitutable Languages}

\begin{lemma}\label{lem:app-u1-insertion}
Let $L\subseteq\Delta^+$ be finite and Clark--Eyraud substitutable.  There are no
$P,Q\in\Delta^*$ and $T\in\Delta^+$ with $PQ\in L$ and $PTQ\in L$.
\end{lemma}

\begin{proof}
Clark--Eyraud substitutability is $(1,s)$-tuple substitutability for the trivial
typing.  Apply Lemma~\ref{lem:no-insertion} to $x=P$, $y=PT$ and the context
$\blank_1Q$, or, if $P=\eps$, to $x=Q$, $y=TQ$ and the context $\blank_1$.
\end{proof}

\begin{lemma}\label{lem:app-u1-code}
Let $h:\Delta^*\to\Sigma^*$ be a morphism such that every $h(\delta)$,
$\delta\in\Delta$, is nonempty and no letter occurs twice in the concatenation of
all $h(\delta)$.  If $L\subseteq\Delta^+$ is finite and Clark--Eyraud
substitutable, so is $h(L)$.
\end{lemma}

\begin{proof}
Each letter occurring in some $h(\delta)$ determines $\delta$ and its position in
$h(\delta)$.  Hence $h$ is injective, and in a word $h(D)$ every letter occurrence
belongs to a determined occurrence of a codeword $h(\delta)$; whether two adjacent
letters lie in one codeword occurrence, and whether a letter begins or ends one,
is read off the letters themselves.

Let $x,y\in\Sigma^+$ with $pxq=h(D_1)$ and $pyq=h(D_2)$ in $h(L)$, $x\ne y$.  The
codeword occurrences lying inside $p$, and those inside $q$, are the same in both
words; so $D_i=PM_iQ$, where $M_i$ consists of the codeword occurrences meeting
$x$, respectively $y$.  If $p$ ends inside a codeword occurrence $h(\delta)$, then
$M_1$ and $M_2$ both begin with $\delta$, and if $q$ begins inside one,
$h(\delta')$, both end with $\delta'$.

Suppose that $x$ lies inside one codeword occurrence $h(\delta)$ and is not all
of it; say $p$ ends inside it (the other case is symmetric).  Then $M_1=\delta$
and $M_2$ begins with $\delta$; $M_2\ne M_1$, since otherwise $y=x$, so
$M_2=\delta T$ with $T\in\Delta^+$, and $D_1=P\delta Q$, $D_2=P\delta TQ$ contradict
Lemma~\ref{lem:app-u1-insertion}.  The
same holds with $x$ and $y$ exchanged.  So neither lies inside one codeword
occurrence unless it is a whole codeword.  Then $M_1=[\delta]X[\delta']$ and
$M_2=[\delta]Y[\delta']$, where $X,Y\in\Delta^*$ are the codewords lying inside $x$,
respectively $y$, and the bracketed letters are present iff the corresponding
cut lies inside a codeword.  If $X=\eps\ne Y$ or $Y=\eps\ne X$,
Lemma~\ref{lem:app-u1-insertion} is contradicted again, and $X=Y=\eps$ gives
$x=y$.  Hence $X,Y\in\Delta^+$, and
\[
 x=\alpha\,h(X)\,\beta,\qquad y=\alpha\,h(Y)\,\beta,
\]
where $\alpha$ is the part of $h(\delta)$ after $p$ (or $\eps$) and $\beta$ the
part of $h(\delta')$ before $q$ (or $\eps$).  The words $X$ and $Y$ share a context
in $L$, so $\D_L(X)=\D_L(Y)$.

Now let $p'xq'\in h(L)$.  The cut between $\alpha$ and $h(X)$, and the one between
$h(X)$ and $\beta$, are boundaries of codeword occurrences, which is read off the
letters of $x$; so $p'\alpha=h(P')$ and $\beta q'=h(Q')$ for some $P',Q'$, and
$P'XQ'\in L$.  Then $P'YQ'\in L$, and $h(P'YQ')=p'yq'$.  By symmetry
$\D_{h(L)}(x)=\D_{h(L)}(y)$.
\end{proof}


\subsection{Quadratic Data for the Binary Typing}

We prove Proposition~\ref{prop:binary-u1-quadratic}.  Put
$k:=2^m$, $M:=2^n$ and $\psi(w):=|w|_{\mathtt b}+|w|_{\mathtt a}/k$, which is
additive.  For a word $w$ let $\lambda(w)$ and $\varrho(w)$ be the lengths of its
longest suffix and longest prefix in $\mathtt a^*$.

\begin{proof}[Proof of Proposition~\ref{prop:binary-u1-quadratic}]
(i) A word of $\Sigma^+$ lies in $L_{m,n}$ iff all its maximal runs of $\mathtt a$
have length divisible by $k$ and $\psi(w)=M$: such a word factors uniquely into
blocks $\mathtt a^k$ and $\mathtt b$, and $\psi$ counts the blocks.  Let
$x,y\in\Sigma^+$ have the same type and let $uxv,uyv\in L_{m,n}$.  By additivity
$\psi(x)=\psi(y)$.  If $x,y$ are pure, this gives $x=y$.  Otherwise write
$x=\mathtt a^{\alpha}x'\mathtt a^{\beta}$ with $x'$ beginning and ending with
$\mathtt b$, and similarly $y=\mathtt a^{\alpha'}y'\mathtt a^{\beta'}$.  The maximal
$\mathtt a$-runs inside $x'$ and $y'$ are maximal runs of $uxv$ and $uyv$, so their
lengths are divisible by $k$; the runs through the two ends of $x$ have lengths
$\lambda(u)+\alpha$ and $\beta+\varrho(v)$, so $\alpha\equiv\alpha'$ and
$\beta\equiv\beta'\pmod k$.  For any context $(p,q)$ the maximal $\mathtt a$-runs of
$pxq$ are those of $p$ not at its end, a run of length $\lambda(p)+\alpha$, the
runs inside $x'$, a run of length $\beta+\varrho(q)$, and those of $q$ not at its
start; the same holds for $pyq$ with $\alpha',\beta',y'$.  So the divisibility
conditions agree, $\psi(pxq)=\psi(pyq)$, and $pxq\in L_{m,n}\iff pyq\in L_{m,n}$.
Finite languages are context-free.

(ii) Take nonterminals $C_0,\ldots,C_m$, $A_0,\ldots,A_n$, start symbol $A_n$, and the
rules $C_0\to\mathtt a$, $C_j\to C_{j-1}C_{j-1}$ ($1\le j\le m$),
$A_0\to C_m\mid\mathtt b$ and $A_i\to A_{i-1}A_{i-1}$ ($1\le i\le n$).  Then
$L(C_j)=\{\mathtt a^{2^j}\}$, $L(A_0)=\{\mathtt a^k,\mathtt b\}$ and $L(A_i)$ is the
set of concatenations of $2^i$ words of $L(A_0)$, so $L(A_n)=L_{m,n}$; the grammar
is good and reduced.  It has $m+n+2$ nonterminals and rules of sizes $2$, $5$
($m$ rules), $3$, $2$ and $5$ ($n$ rules) (Definition~\ref{def:standard-mcfg}),
which gives $6m+6n+9$.  At fan-out one, the shortest sentence using a rule
$R:B\to\rho(B_1,\ldots,B_q)$ has length $\theta_G(R)=o_B+c_R+\sum_jd_{B_j}$, where
$d_C$ is the length of a shortest yield of $C$, $o_C$ the least number of letters
around an occurrence of $C$ in a sentence, and $c_R$ the number of terminals
written by $R$; here
$d_{A_0}=1$, $d_{A_i}=2^i$, $d_{C_j}=2^j$, $o_{A_i}=M-2^i$, $o_{A_0}=o_{C_m}=M-1$ and
$o_{C_j}=M-1+k-2^j$.  This gives $\theta_G(R)=M$ for the rules of $A_i$ and for
$A_0\to\mathtt b$, and $\theta_G(R)=M+k-1$ for $A_0\to C_m$ and the rules of the
$C_j$; so $\bar\theta_{G_{m,n}}=2^m+2^n-1$.

(iii) The only pure word of $L_{m,n}$ is $\mathtt a^{kM}$, so by
Proposition~\ref{prop:rare-type} every characteristic sample contains it and
$\|F\|_+\ge kM+1>2^{m+n}$.  For $m=n$, $\bar\theta_{G_{n,n}}<2^{n+1}$ gives
$\bar\theta_{G_{n,n}}^{\,2}/4<4^n<\|F\|_+$.  A bound $p(\|G\|)\bar\theta_G$ would give
$4^n\le p(12n+9)\,2^{n+1}$ for all $n$, which fails for every polynomial $p$.
\end{proof}


%%%%%%%%%%%%%%%%%%%%%%%%%%%%%%%%%%%%%%%%%%%%%%%%%%%%%%%%%%%%%%%%%%%%%%
%% REMOVED: apps two-letter, singleton, comparison
%%%%%%%%%%%%%%%%%%%%%%%%%%%%%%%%%%%%%%%%%%%%%%%%%%%%%%%%%%%%%%%%%%%%%%
\section{Content Typings on Two Letters at Fan-Out at Least Two}
\label{app:two-letter-content}

This appendix proves Theorem~\ref{thm:two-letter-content}.  Throughout,
$\Sigma=\{\mathtt a,\mathtt b\}$, $c:=c_\Sigma$, $f\ge2$ is fixed, $L\in\Cmcf{f}{c}$, and
$G=(V,\Sigma,P,S,\fan)$ is a reduced good grammar of fan-out $\le f$ and rank
$\le r:=r_{\max}$ with $L(G)=L\cap\Sigma^+$; put $v:=|V|$ and
$\theta:=\bar\theta_G$; we also put $D_f:=(6f+1)3^f$.  Apart from Lemma~\ref{lem:app-two-basic}(i), which uses the
condition of the arity of the typed nonterminal at hand, the arguments use only
$(1,c)$- and $(2,c)$-tuple-substitutability, which hold since $f\ge2$.  For each $A\in V$ the context $E_A$ is that of
Lemma~\ref{lem:thickness-equivalence}(iv), so $|E_A|_{\mathrm{fix}}\le\theta$
and $E_A[L_A(G)]\subseteq L(G)$; for each $B\in V$ fix $\omega_B\in L_B(G)$ with
$\|\omega_B\|_1\le\theta$ (item~(ii)), and recall that every rule writes at most
$\theta$ terminals (item~(i)).  In a derivation tree, the \emph{yield} of a node
$u$ is the tuple it derives, $\ell_u$ is its total length, $B_u$ its
nonterminal, and $u$ is \emph{heavy} if $\ell_u>\theta$; a node \emph{writes}
the terminals of its template.  The letters $x,y$ always denote the two
letters of $\Sigma$ in some order.  A \emph{run} of a word is a maximal factor
$x^k$ ($k\ge1$); a run is \emph{internal} if it is neither the first nor the
last run, so that an internal $y$-run is preceded and followed by $x$.  The
\emph{$x$-skeleton} $\sigma_x(w)$ replaces every $x$-run of $w$ by one $x$, and
$w$ has a \emph{spare} $x$ if $|w|_x$ exceeds the number of $x$-runs of $w$.
We say that \emph{$x^k$ occurs} if some word of $L$ has the factor $x^k$, and
that $L$ is \emph{Parikh-saturated} if for $u,w\in\Sigma^+$ with the same
Parikh vector $\Psi(u)=\Psi(w)$ we have $u\in L$ iff $w\in L$.

\paragraph{Map of the proof.}
The table lists, for each step, what it uses and where it is used.
\begin{center}
\footnotesize
\setlength{\tabcolsep}{3pt}
\renewcommand{\arraystretch}{1.15}
\begin{tabularx}{\linewidth}{@{}>{\raggedright\arraybackslash}p{0.24\linewidth}>{\raggedright\arraybackslash}X>{\raggedright\arraybackslash}p{0.26\linewidth}@{}}
\hline
Step & Uses & Used in\\
\hline
Lemma~\ref{lem:app-two-basic}(i) typed closure
 & Lemma~\ref{lem:typed-homogeneity}, arity-$d$ substitutability
 & Lemma~\ref{lem:app-two-saturated}\\
Lemma~\ref{lem:app-two-basic}(ii) cutting identity paths
 & goodness of $G$
 & (P2), the mixed case of Lemma~\ref{lem:app-two-saturated}\\
Lemma~\ref{lem:app-two-runs} run exchanges
 & $(1,c)$, $(2,c)$-substitutability, $x^4$ occurs
 & Lemma~\ref{lem:app-two-alternate}, Proposition~\ref{prop:app-two-nonsaturated}\\
Lemma~\ref{lem:app-two-alternate}
 & Lemma~\ref{lem:app-two-runs}(i),(iii)
 & Proposition~\ref{prop:app-two-nonsaturated}\\
Lemma~\ref{lem:app-two-insertion} runs created by replacement
 & Lemma~\ref{lem:witness-length-template}(i), node-minimal shortest $T$
 & Proposition~\ref{prop:app-two-nonsaturated}\\
Proposition~\ref{prop:app-two-nonsaturated} (not saturated)
 & the three lemmas above; typed labels on paths
 & Theorem~\ref{thm:two-letter-content}\\
Lemma~\ref{lem:app-two-saturated}, union $\Sigma$
 & Lemma~\ref{lem:app-two-basic}; floors along paths
 & Theorem~\ref{thm:two-letter-content}\\
(P1) row rigidity
 & $(1,c)$-substitutability, minimality of $K$
 & (P3), (P4)\\
(P2) height, mixed nodes
 & Lemma~\ref{lem:app-two-basic}(ii), minimality of $K$
 & (P3), (P6)\\
(P3) heavy pure nodes
 & (P1), (P2)
 & (P4), (P5)\\
(P4) common ratio $\alpha$
 & (P3)
 & (P5), (P6)\\
(P5) bound on $\alpha$
 & (P3), (P4)
 & (P6)\\
(P6) counting
 & (P2), (P4), (P5)
 & Theorem~\ref{thm:two-letter-content}\\
\hline
\end{tabularx}
\end{center}

\begin{lemma}\label{lem:app-two-basic}
\textup{(i)} For every typed nonterminal $X=A_{\mathbf t}$ of arity $d$,
$\operatorname{Cl}_{L,c}(X)=\{\tuple z\in(\Sigma^+)^d:c^{(d)}(\tuple z)=\mathbf t,\
E_A[\tuple z]\in L\}$; hence $\tau_X$ is the least total length of such a tuple.

\textup{(ii)} In a derivation tree of $G$, let $u$ be a proper ancestor of $u'$
with the same nonterminal, and suppose their yields have the same total Parikh
vector.  Then every rule on the path from $u$ to $u'$ has rank one, writes no
terminal and is a componentwise identity, so replacing the subtree at $u$ by the
subtree at $u'$ changes no derived tuple and removes at least one node.
\end{lemma}

\begin{proof}
(i) Fix $\tuple u\in L_X$.  By Lemma~\ref{lem:typed-homogeneity},
$\operatorname{Cl}_{L,c}(X)$ is the set of tuples with the type and the
distribution of $\tuple u$; they are accepted by $E_A\in\D_L(\tuple u)$.  A tuple
$\tuple z$ of type $\mathbf t$ accepted by $E_A$ shares this context with $\tuple u$,
so $\D_L(\tuple z)=\D_L(\tuple u)$ by substitutability.

(ii) By linearity and nondeletion every letter of the yield of $u'$ occurs in the
yield of $u$, so equal Parikh vectors exclude terminals written on the path
and siblings off the path, since siblings have nonempty yields
($\lambda$-freeness).  So every rule on the path has one child of some fan-out
$d'$, each of its $d'$ variables occurs once (nondeletion) and in order
(nonpermutation), and each parent component is nonempty ($\lambda$-freeness).
Two variables in one component would be consecutive, which nonmerging forbids;
hence every component is a single variable, the fan-outs agree, and the rule is
the identity.
\end{proof}

\begin{lemma}[Run exchanges]\label{lem:app-two-runs}
Suppose that $x^4$ occurs.
\begin{enumerate}[label=\textup{(\roman*)},leftmargin=*]
\item If $w,w'\in\Sigma^+$ satisfy $\sigma_x(w)=\sigma_x(w')$ and
$|w|_x=|w'|_x$, then $w\in L$ iff $w'\in L$.
\item If some word of $L$ has a factor $xxyx$ or $xyxx$, then $L$ is
Parikh-saturated.
\item If $L$ is not Parikh-saturated, then in every word of $L$ with a factor
$xx$ every internal $y$-run has length at least $2$.
\item If $L$ is not Parikh-saturated and $y^4$ does not occur, then every word
of $L$ with a factor $xx$ has at most one internal $y$-run.
\end{enumerate}
\end{lemma}

\begin{proof}
(i) Since $c(x)$ is idempotent, Lemma~\ref{lem:power-transfer} applies with
$u=x$ and $a=b=d=1$: as $x^4$ occurs, $w_{\mathrm L}xw_{\mathrm M}xxw_{\mathrm R}\in L$ iff
$w_{\mathrm L}xxw_{\mathrm M}xw_{\mathrm R}\in L$ whenever $w_{\mathrm M}\ne\eps$.  Let $R$ and $R'$ be $x$-runs of a word $w$ with
$|R|\ge2$, say $R'$ to the left of $R$ (the other case is symmetric).  Write
$w=U\,x\,Z\,xx\,V$, where the first displayed $x$ is the last letter of $R'$ and
$xx$ are the first two letters of $R$; the factor $Z$ separates two runs, so
$Z\ne\eps$.  Hence $U\,xx\,Z\,x\,V\in L$ iff $w\in L$.  This moves one $x$ from $R$
to $R'$ and keeps the $x$-skeleton and $|w|_x$.  Words with the same
$x$-skeleton have the same number $\rho$ of $x$-runs, and their run lengths are
compositions of $|w|_x$ into $\rho$ positive parts; any two such compositions
are connected by moves of one unit from a part of size at least $2$ to another
part.

(ii) Say $UxxyxV\in L$ (the other case is symmetric).  By (i),
$UxyxxV\in L$.  The context $(Ux,xV)$ accepts $xy$ and $yx$, both of content
$\Sigma$, so $\D_L(xy)=\D_L(yx)$: in every context $xy$ may be replaced by $yx$
and conversely.  Adjacent transpositions connect all words with one Parikh
vector.

(iii) Let $w\in L$ have a factor $xx$ and an internal $y$-run of length $1$.
Then $w$ has a spare $x$, and by (i) the spare letters can be moved into the
$x$-run preceding that $y$-run; the resulting word of $L$ has a factor $xxyx$,
so $L$ is Parikh-saturated by (ii).

(iv) Suppose that $w\in L$ has a factor $xx$ and two internal $y$-runs.  Then it
has two consecutive internal $y$-runs $y^e$ and $y^{e'}$, separated by an
$x$-run $R$; let $A_0$ be the $x$-run preceding $y^e$.  By (iii), $e,e'\ge2$.
Since $w$ has a spare $x$, (i) gives two words of $L$ with the $x$-skeleton and
the number of letters $x$ of $w$: a word $W_2=U\,x\,y^e\,xx\,V'$ in which $R$ has
length $2$ (the displayed $x$ before $y^e$ is the last letter of $A_0$, and $V'$
begins with $y^{e'}$), and a word $W_1=U_1\,x\,y^e\,x\,y^{e'}V_1$ in which $R$ has
length $1$.  The context $(U,xV')$ accepts $xy^ex$, giving $W_2$, and $xxy^e$,
giving $U\,xx\,y^e\,x\,V'$, which has the $x$-skeleton and the $x$-count of $W_2$
and lies in $L$ by (i).  Both tuples have content $\Sigma$, so
$\D_L(xy^ex)=\D_L(xxy^e)$.  The context $(U_1,y^{e'}V_1)$ accepts $xy^ex$, giving
$W_1$, hence also $xxy^e$, so $U_1\,xx\,y^{e+e'}V_1\in L$.  As $e+e'\ge4$, this
word has the factor $y^4$, a contradiction.
\end{proof}

\begin{lemma}\label{lem:app-two-alternate}
If $L$ is not Parikh-saturated and both $\mathtt a^4$ and $\mathtt b^4$ occur,
then every run of every word of $L$ with at least four runs has length $1$.
\end{lemma}

\begin{proof}
Let $w\in L$ have at least four runs.  Its second and third runs are internal
and consist of different letters, so for each letter $y$ the word $w$ has an
internal $y$-run $\beta$ and at least two $y$-runs.  By
Lemma~\ref{lem:app-two-runs}(i) for $y$, moving letters $y$ from $\beta$ to another
$y$-run gives a word of $L$ in which $\beta$ has length $1$ and the $x$-runs are
those of $w$.  If an $x$-run of $w$ had length at least $2$, this word would
have a factor $xx$ and an internal $y$-run of length $1$, contrary to
Lemma~\ref{lem:app-two-runs}(iii).  So all $x$-runs of $w$ have length $1$, for
both choices of $x$.
\end{proof}

\begin{lemma}[Runs created by short replacements]\label{lem:app-two-insertion}
Let $X=A_{\mathbf t}$ be a typed nonterminal of arity $d\le f$, and let $T$ be a
derivation of $\widetilde G_0$ from $X$ whose yield $\tuple u$ has the least total
length and, among those, the fewest nodes.  If $T$ has an antichain $Q$ of
heavy nodes with $|Q|\ge6f+2$, then for some letter $x$:
\textup{(a)} some component of $\tuple u$ has content $\{x\}$ and length at least
$5$, so $x^5$ occurs; and
\textup{(b)} some word $W\in L$ has a factor $xx$ and at least four runs, and,
unless both $\mathtt a^5$ and $\mathtt b^5$ occur, $W$ has two internal $y$-runs.
\end{lemma}

\begin{proof}
In each component $u_i$ of content $\{x\}$ (a \emph{pure} component) mark the
first and the last $\min\{2,|u_i|\}$ letters, and in each component of content
$\Sigma$ mark one $\mathtt a$ and one $\mathtt b$; there are at most $4f$ marks.
Each mark is written by one node, and nodes of $Q$ have disjoint subtrees, so
the set $Q_0$ of nodes of $Q$ whose subtrees write no mark has
$|Q_0|\ge|Q|-4f\ge2f+2$.

Let $h\in Q_0$ with nonterminal $C_h$.  Erase the type indices and replace the
subtree at $h$ by a derivation of $\omega_{C_h}$.  The new root tuple lies in
$L_A(G)$ and is shorter, since by nondeletion every letter of the yield of $h$
occurs in $\tuple u$ and $\|\omega_{C_h}\|_1\le\theta<\ell_h$.  The marks keep the
content of every component of content $\Sigma$.  If every pure component stayed
pure, the new tuple would have type $\mathbf t$ and would lie in $L_X$ by
Lemma~\ref{lem:witness-length-template}(i) with $\lambda:=\widehat c$, contrary to
the choice of $T$.  Hence $\omega_{C_h}$ puts the other letter into a block of $h$
(an interval of a root component occupied by a component of the yield of $h$)
that lies in a pure component.  Blocks contain no mark, so that pure component
has length at least $5$, and the block lies after its first two and before its
last two letters.  This proves (a), since $E_A[\tuple u]\in L$.

Each $h\in Q_0$ owns at most $f$ blocks.  Join two owners when blocks of theirs in
one root component are adjacent with no letter between them.  Each block has at
most two adjacent blocks, so this graph has maximum degree at most $2f$ and an
independent set $I$ with $|I|\ge|Q_0|/(2f+1)>1$, hence $|I|\ge2$.  Replace the subtrees at
all $h\in I$ at once, giving $\tuple u'\in L_A(G)$ and $W:=E_A[\tuple u']\in L$.
Fix a pure component of letter $x$ and write $y$ for the other letter.  Its
letters outside the blocks of $I$ are unchanged, hence equal to $x$; this
covers the blocks of owners outside $I$, which are kept, and the first two and
the last two letters, which are marked.  So every maximal $y$-run of the new
component is preceded and followed by $x$, and it lies in a union of blocks of
$I$ that are consecutive with no letter between them.  Such blocks belong to one
owner: two blocks of distinct owners $h\ne h'$ of $I$ with no block of $I$
between them are separated by an unchanged letter, since their gap is nonempty
(otherwise $h$ and $h'$ would be joined) or contains a block of an owner outside
$I$.  Hence every $y$-run of the new component has a unique \emph{creating}
owner in $I$, and by the previous paragraph every $h\in I$ creates at least one
run, in a pure component whose letter we call $x_h$.  For $h\ne h'$ in $I$ the
runs they create are distinct: if they lie in the same pure component, they have
different creating owners; if they lie in different pure components, they are
in different components of $\tuple u'$.

Take $h_1\ne h_2$ in $I$.  If $x_{h_1}=x_{h_2}=:x$, then $W$ has the factor $xx$
(the first two letters of a pure component of length at least $5$) and two
internal $y$-runs, hence at least five runs.  Otherwise both $\mathtt a^5$ and
$\mathtt b^5$ occur, by the argument for (a) applied to $h_1$ and to $h_2$, and with $x:=x_{h_1}$ and $y:=x_{h_2}$ the two pure
components contain the runs $x,y,x$ and $y,x,y$ respectively; these are six
distinct runs of $W$, and $W$ has the factor $xx$.
\end{proof}

\begin{proposition}\label{prop:app-two-nonsaturated}
If $L$ is not Parikh-saturated, then $t_{\mathrm{typ}}(G,c)\le D_f(r+1)v\theta$.
\end{proposition}

\begin{proof}
Let $X$ and $T$ be as in Lemma~\ref{lem:app-two-insertion}.  We first show that
every antichain of heavy nodes of $T$ has at most $6f+1$ nodes.  Otherwise
Lemma~\ref{lem:app-two-insertion} applies.  If neither $\mathtt a^4$ nor
$\mathtt b^4$ occurs, (a) is impossible.  If both occur, the word $W$ of (b)
has at least four runs and a run of length at least $2$, contrary to
Lemma~\ref{lem:app-two-alternate}.  If exactly one occurs, say $x_0^4$, then by
(a) $x=x_0$ and $y^5$ does not occur, so by (b) $W$ has a factor $xx$ and two
internal $y$-runs, contrary to Lemma~\ref{lem:app-two-runs}(iv).

A typed label is a nonterminal together with the contents of its components,
so there are at most $3^fv$ of them.  No typed label repeats on a path of $T$,
since cutting between two equal labels gives a derivation from $X$ whose yield
is not longer and which has fewer nodes; so paths have at most $3^fv$ nodes.
Heavy nodes are closed under ancestors, and those at one depth form an
antichain, so there are at most $(6f+1)3^fv=D_fv$ of them.  If the root is not
heavy, $\|\tuple u\|_1\le\theta$.  Otherwise every letter of $\tuple u$ is written by
a heavy node or lies in the yield of a non-heavy child of one, and each heavy
node writes at most $\theta$ letters and has at most $r$ children.  Hence
$\|\tuple u\|_1\le D_f(r+1)v\theta$.
\end{proof}

\begin{lemma}\label{lem:app-two-saturated}
If $L$ is Parikh-saturated, then
$\tau(G,c)\le C_{f,r}v^2\theta^4H_\theta$, where $C_{f,r}:=(r+1)(4r(3f+1)+D_f)+3$.
\end{lemma}

\begin{proof}
Fix $X=A_{\mathbf t}$ of arity $d$, put $(p,q):=\Psi(E_A)$, so $p+q\le\theta$, and
$\mathcal P:=\Psi(L\cap\Sigma^+)$.  By Lemma~\ref{lem:app-two-basic}(i) and Parikh
saturation, $\tau_X$ is the least $\|\tuple z\|_1$ over tuples $\tuple z$ of type
$\mathbf t$ with $(p,q)+\sum_i\Psi(z_i)\in\mathcal P$.  Write $(x_B,y_B):=\Psi(\omega_B)$,
and for a node $u$ of a derivation tree let $(a_u,b_u)$ be the Parikh vector of
its yield.  If $u$ lies in a successful derivation of a word $w$, replacing its
subtree by a derivation of $\omega_{B_u}$ gives a word of $L$ with Parikh vector
$\Psi(w)-(a_u,b_u)+(x_{B_u},y_{B_u})$; for a proper descendant $u'$ with
$B_{u'}=B_u$, replacing the subtree at $u$ by that at $u'$ gives one with
$\Psi(w)-(a_u-a_{u'},b_u-b_{u'})$.  Every comparison below is of this kind: it
concerns Parikh vectors of words of $L$ yielded by complete successful
derivations obtained from a fixed one by such replacements, and by Parikh
saturation membership in $L$ of a word depends only on its Parikh vector.  The union of the contents $t_1,\dots,t_d$ is
$\Sigma$, or all of them equal one singleton.

\emph{The union is $\Sigma$.}  Let $t_{\mathtt a}$ and $t_{\mathtt b}$ be the
numbers of components whose content contains $\mathtt a$, resp.\ $\mathtt b$; both
lie in $\{1,\dots,f\}$.  A tuple of type $\mathbf t$ with $n_{\mathtt a}$ letters $\mathtt a$
and $n_{\mathtt b}$ letters $\mathtt b$ exists iff $n_{\mathtt a}\ge t_{\mathtt a}$ and
$n_{\mathtt b}\ge t_{\mathtt b}$, so $\tau_X$ is the least
$n_{\mathtt a}+n_{\mathtt b}$ with these bounds and
$(p+n_{\mathtt a},q+n_{\mathtt b})\in\mathcal P$.  Fix a minimizing pair, put
$c_{\mathtt a}:=n_{\mathtt a}-t_{\mathtt a}+1\ge1$ and
$c_{\mathtt b}:=n_{\mathtt b}-t_{\mathtt b}+1\ge1$, and fix a successful derivation
$T$ of $W:=\mathtt a^{p+n_{\mathtt a}}\mathtt b^{q+n_{\mathtt b}}$ with the fewest
nodes.  If a heavy node $u$ had $a_u<c_{\mathtt a}$ and $b_u<c_{\mathtt b}$, its
replacement would give a Parikh vector $(p+n'_{\mathtt a},q+n'_{\mathtt b})\in\mathcal
P$ with $n'_{\mathtt a}\ge t_{\mathtt a}$, $n'_{\mathtt b}\ge t_{\mathtt b}$ and
$n'_{\mathtt a}+n'_{\mathtt b}<n_{\mathtt a}+n_{\mathtt b}$; so every heavy node has
$a_u\ge c_{\mathtt a}$ or $b_u\ge c_{\mathtt b}$.  If $u'$ is a proper descendant of
$u$ with the same nonterminal, then $a_u-a_{u'}\ge c_{\mathtt a}$ or
$b_u-b_{u'}\ge c_{\mathtt b}$: a zero difference contradicts node-minimality
(Lemma~\ref{lem:app-two-basic}(ii)), and another small difference contradicts
the choice of the pair.  Along a path, $(\lfloor a_u/c_{\mathtt a}\rfloor,
\lfloor b_u/c_{\mathtt b}\rfloor)$ is componentwise nonincreasing, and no
nonterminal repeats while it is constant.  Since
$\lfloor(p+n)/(n-t+1)\rfloor\le p+t$ for $n\ge t\ge1$ and $p\ge0$ (as $(p+t)(n-t+1)-(p+n)=(n-t)(p+t-1)\ge0$), it
takes at most $1+(p+t_{\mathtt a})+(q+t_{\mathtt b})\le\theta+2f+1$ values, so paths
have at most $v(\theta+2f+1)$ nodes.  Heavy nodes at one depth have disjoint yields,
so there are at most $(p+n_{\mathtt a})/c_{\mathtt a}+(q+n_{\mathtt b})/c_{\mathtt
b}\le\theta+2f$ of them.  Heavy nodes are closed under ancestors, and every letter
is written by a heavy node or lies in one of its at most $r$ non-heavy
children; hence
$\tau_X\le|W|\le\max\{\theta,(r+1)\theta(\theta+2f)(\theta+2f+1)v\}\le(2f+1)(2f+2)(r+1)v\theta^3$.

\emph{All contents equal $\{\mathtt a\}$} (the case $\{\mathtt b\}$ is symmetric),
and $h:=d\in\{1,\dots,f\}$.  Let $K:=\min\{k\ge h:(p+k,q)\in\mathcal P\}$; then
$\tau_X=K$, and we may assume $K>\theta+2(f-1)\ge p+2(h-1)$, since otherwise
$\tau_X\le(2f-1)\theta$.  Fix a successful derivation
$T$ of $W:=\mathtt b^q\mathtt a^{p+K}$ with the fewest nodes.

(P1) \emph{Row rigidity}: there are no $(c_0,y),(c_0+\delta,y)\in\mathcal P$ with
$c_0\ge1$ and $1\le\delta\le K-h$.  Otherwise the context
$(\mathtt b^y\mathtt a^{c_0-1},\eps)$ accepts $\mathtt a$ and $\mathtt a^{1+\delta}$, so
$\D_L(\mathtt a)=\D_L(\mathtt a^{1+\delta})$ by $(1,c)$-substitutability (both have
content $\{\mathtt a\}$); the context $(\mathtt b^q\mathtt a^{p+K-\delta-1},\eps)$ accepts
$\mathtt a^{1+\delta}$, hence $\mathtt a$, and $(p+K-\delta,q)\in\mathcal P$ with
$h\le K-\delta<K$, contradicting the choice of $K$.

(P2) \emph{Height.}  If $u'$ is a proper descendant of $u$ with the same
nonterminal and $b_u=b_{u'}$, then $a_u-a_{u'}\ge K-h+1$: the value $0$
contradicts node-minimality (Lemma~\ref{lem:app-two-basic}(ii)), and a value in
$[1,K-h]$ contradicts the choice of $K$.  Going down a path, $b_u$ is
nonincreasing (by nondeletion the yield of a child occurs in that of its parent),
so the nodes of the path with one value of $b_u$ form a segment; $b_u$ takes at most $q+1$
values, and three nodes with one nonterminal and one value of $b_u$ would give
$p+K\ge a_u\ge2(K-h+1)$, false as $K>p+2h-2$.  So paths have at most $2v(q+1)$
nodes.  A node is \emph{mixed} if its yield contains $\mathtt b$.  For
$1\le k\le q$, the mixed nodes with $b_u=k$ that have no such proper ancestor have
pairwise disjoint yields, so there are at most $\lfloor q/k\rfloor$ of them; below
one of them, two such nodes are comparable (otherwise its yield would contain $2k$
letters $\mathtt b$), so they form a chain, which has at most $2v$ nodes by the
argument above.  Hence there are
$M\le2v\sum_{k=1}^q\lfloor q/k\rfloor\le2vqH_q\le2v\theta H_\theta$ mixed nodes.

(P3) \emph{Heavy pure nodes.}  Call a node $u$ that is not mixed a
\emph{heavy pure} node if $\ell_u=a_u>\theta+f-1$, and put
$g_u:=\ell_u-x_{B_u}\ge f$.  Two comparable heavy pure nodes have different
nonterminals, since by (P2) the upper one would have
$\ell_u\ge K-h+1+\theta+f>p+K=|W|_{\mathtt a}$.  In particular heavy pure nodes with
one nonterminal have disjoint yields.  For an antichain $J$ of heavy pure nodes,
replacing all of them gives a word with Parikh vector $(p+K-\sigma(J),q+Y(J))$,
where $\sigma(J):=\sum_{u\in J}g_u$ and $Y(J):=\sum_{u\in J}y_{B_u}$; this word is yielded by a complete successful derivation, so its Parikh vector
lies in $\mathcal P$.  Hence two antichains $J,J'$ with
$\sigma(J),\sigma(J')\le K-h$ and $Y(J)=Y(J')=:Y$ give
$(p+K-\sigma(J),q+Y),(p+K-\sigma(J'),q+Y)\in\mathcal P$, whose first coordinates
are at least $p+h\ge1$ and differ by $|\sigma(J)-\sigma(J')|\le K-h$; by (P1),
$\sigma(J)=\sigma(J')$.

(P4) \emph{Repeated nonterminals.}  If $B$ labels $m_B\ge2$ heavy pure nodes,
their yields are disjoint by (P3), so these nodes are pairwise incomparable and
every set of them is an antichain, so each has $\ell_u\le p+K-(\theta+f)\le K-f$ and
$g_u\le K-f\le K-h$; comparing two singletons by (P3) gives a common value $g_B$,
and then a common $\ell_B$.  If moreover $y_B=0$, one replacement gives
$(p+K-g_B,q)\in\mathcal P$ with $h\le K-g_B<K$, which is impossible; so $y_B\ge1$.
Call $B$ \emph{frequent} if $m_B>\theta$.  For frequent $B,C$, choose $y_C$ nodes
labelled $B$ and, separately, $y_B$ nodes labelled $C$, which is possible since
$m_B>\theta\ge y_C$ and $m_C>\theta\ge y_B$; both sets are antichains with
$Y=y_By_C$, and $y_Cg_B\le(m_B-1)\ell_B\le p+K-\ell_B\le K-f\le K-h$, and likewise
for $C$.  By (P3), $y_Cg_B=y_Bg_C$, so all frequent $B$ share a ratio
$\alpha:=g_B/y_B$, and $\ell_B=x_B+\alpha y_B\le(1+\alpha)\theta$.

(P5) \emph{The ratio.}  Put $\alpha_*:=(r+1)(\theta+f-1)(2\theta+1)v\le3f(r+1)v\theta^2$;
we show $\alpha\le\alpha_*$ when frequent nonterminals exist.  This is clear if
$\alpha\le\theta$.  Otherwise fix a heavy pure node $u$ labelled by a frequent
nonterminal $B$, and let $H_u$ be the number of heavy pure nodes in its subtree.
These nodes are closed under ancestors within that subtree.  Every letter in the
yield of $u$ is written by a heavy pure node or belongs to the yield of a
non-heavy child of such a node.  Each heavy pure node writes at most $\theta$
letters and has at most $r$ such children, each of total yield length at most
$\theta+f-1$; hence $\ell_B\le(r+1)(\theta+f-1)H_u$.  Group these nodes by their
number of heavy pure proper ancestors in the subtree.  A chain of heavy pure
nodes has pairwise different nonterminals by (P3), hence at most $v$ nodes, so
this number is less than $v$ and there are at most $v$ groups; two comparable
nodes have different numbers, so each group is an antichain.  In one group the nodes with frequent labels have
length at least $\alpha$ and disjoint yields inside that of $u$, so there are at
most $(1+\alpha)\theta/\alpha<\theta+1$ of them; nodes with other labels number at
most $v\theta$ in all of $T$.  So $H_u\le v(2\theta+1)$, and
$\alpha\le g_B\le\ell_B\le\alpha_*$.

(P6) \emph{Counting.}  Call a node that is not mixed \emph{large} if
$\ell_u>T_*:=(\alpha_*+1)\theta\le(3f+1)(r+1)v\theta^3$.  By (P4)--(P5) no frequent
nonterminal labels a large node, so there are at most $v\theta$ large nodes, and
they are closed under ancestors among nodes that are not mixed.  Mixed nodes are closed under ancestors, since by
nondeletion the yield of a node occurs in that of its parent.  Let a letter
$\mathtt a$ of $W$ be written by a node $w$ that is neither mixed nor large.  By the
two closure properties, the ancestors of $w$ (including $w$) that are neither
mixed nor large form a final segment of the path from the root to $w$; let $z$ be
its top node.  Then $z$ is the root of $T$ or a child of a mixed or a large node,
and the subtree at $z$ contains no mixed or large node, so $\ell_z\le T_*$.  Thus
every letter $\mathtt a$ of $W$ is written by a mixed node, is written by a large
node, or lies in the yield of such a $z$; there are at most $1+rM+rv\theta$ nodes
$z$.  Hence
\[
\begin{aligned}
 p+K&\le\theta M+(rM+rv\theta+1)T_*+\theta\cdot v\theta\\
 &\le3v\theta^2H_\theta+4rv\theta H_\theta\cdot(3f+1)(r+1)v\theta^3
 \le\bigl(4r(3f+1)(r+1)+3\bigr)v^2\theta^4H_\theta,
\end{aligned}
\]
using $M\le2v\theta H_\theta$ and $rM+rv\theta+1\le4rv\theta H_\theta$.  As
$(2f+1)(2f+2)\le D_f$, in both cases, and when
$K\le\theta+2(f-1)$, $\tau_X\le C_{f,r}v^2\theta^4H_\theta$.
\end{proof}

\begin{proof}[Proof of Theorem~\ref{thm:two-letter-content}]
If $L$ is Parikh-saturated, Lemma~\ref{lem:app-two-saturated} applies.  Otherwise
Proposition~\ref{prop:app-two-nonsaturated} gives
$\tau(G,c)\le t_{\mathrm{typ}}(G,c)\le D_f(r+1)v\theta$.  In both cases
$\tau(G,c)\le C_{f,r}v^2\theta^4H_\theta$.  The image of $c$ has three elements, so
$|\mathsf S^\sharp|=4$ and $N_{\mathrm{typ}}\le C_c(G)\le4^{fr}\|G\|$.  By
Corollary~\ref{cor:distribution-representatives} there is a common characteristic
sample with
$\|K\|_+\le N_{\mathrm{typ}}(1+2\theta+r\tau(G,c))+\mathbf 1_{\{\eps\in L\}}$.  With
$v\le\|G\|$, this gives the first bound using $1+2\theta\le3\|G\|^2\theta^4H_\theta$ and
$\tau\le C_{f,r}\|G\|^2\theta^4H_\theta$, and the second using $1+2\theta\le3\|G\|\theta$ and
$\tau\le D_f(r+1)\|G\|\theta$.
\end{proof}

\section{Single Words with Two Holes}
\label{app:singleton}

This appendix proves the classification used in
Proposition~\ref{prop:unique-two-letter-shapes}.  Throughout the typing is
trivial, $W\in\{\mathtt a,\mathtt b\}^+$, and \emph{$W$ passes for arity $d$} if
$\{W\}$ satisfies the condition of Definition~\ref{def:tuple-substitutable} for
that $d$; that is, whenever $E[\tuple x]=E[\tuple y]=F[\tuple x]=W$ for
$E,F\in\mathcal E_d$ and $\tuple x,\tuple y\in(\Sigma^+)^d$, also $F[\tuple y]=W$.
Every word passes for arity $1$, since $pxq=pyq$ forces $x=y$.  A quadruple
$(E,\tuple x,\tuple y,F)$ with $E[\tuple x]=E[\tuple y]=F[\tuple x]=W\ne F[\tuple y]$
is a \emph{violation}.  Exchanging the letters and reversing words map violations
to violations, and a violation for a factor of $W$ gives one for $W$
(Lemma~\ref{lem:unique-word-rigid}).

\begin{theorem}\label{thm:app-singleton}
Let $W\in\{\mathtt a,\mathtt b\}^+$ with $|W|=n\ge7$.  The following are equivalent:
\textup{(i)} $W$ passes for arity $2$; \textup{(ii)} $W$ passes for every arity;
\textup{(iii)} for some $\{x,y\}=\{\mathtt a,\mathtt b\}$, $W$ is $x^n$, $x^{n-1}y$,
$yx^{n-1}$, $yx^{n-2}y$, or alternating.
\end{theorem}

\paragraph{Eleven violations.}
Table~\ref{tab:singleton-certificates} gives a violation for arity $2$ for each of
eleven words of length $5$ to $7$; contexts are written $(v_0,v_1,v_2)$ for
$v_0\blank_1v_1\blank_2v_2$, and each row is checked by concatenation.  Write
$u\sqsubseteq W$ if $W$ has a factor that is $u$ or its image under exchange of
letters, reversal, or both.  By the remarks above, $W$ fails for arity $2$
whenever $u\sqsubseteq W$ for a word $u$ of the table.

\begin{table}[htbp]
\caption{Violations for arity $2$: $E[\tuple x]=E[\tuple y]=F[\tuple x]=W\ne F[\tuple y]$.}
\label{tab:singleton-certificates}
\centering
\scriptsize
\setlength{\tabcolsep}{3pt}
\renewcommand{\arraystretch}{1.05}
\begin{tabular}{@{}llllll@{}}
\hline
$W$ & $E$ & $\tuple x$ & $\tuple y$ & $F$ & $F[\tuple y]$\\
\hline
$\mathtt{aaabb}$ & $(\eps,\mathtt{a},\mathtt{b})$ & $(\mathtt{aa},\mathtt{b})$ & $(\mathtt{a},\mathtt{ab})$ & $(\eps,\mathtt{ab},\eps)$ & $\mathtt{aabab}$\\
$\mathtt{aabba}$ & $(\mathtt{a},\mathtt{b},\eps)$ & $(\mathtt{a},\mathtt{ba})$ & $(\mathtt{ab},\mathtt{a})$ & $(\eps,\mathtt{ab},\eps)$ & $\mathtt{ababa}$\\
$\mathtt{aaaaba}$ & $(\eps,\mathtt{a},\mathtt{ba})$ & $(\mathtt{aa},\mathtt{a})$ & $(\mathtt{a},\mathtt{aa})$ & $(\eps,\mathtt{aab},\eps)$ & $\mathtt{aaabaa}$\\
$\mathtt{aaabaa}$ & $(\mathtt{a},\mathtt{a},\eps)$ & $(\mathtt{a},\mathtt{baa})$ & $(\mathtt{aab},\mathtt{a})$ & $(\eps,\mathtt{aa},\eps)$ & $\mathtt{aabaaa}$\\
$\mathtt{aaabab}$ & $(\eps,\mathtt{a},\mathtt{ab})$ & $(\mathtt{a},\mathtt{ab})$ & $(\mathtt{aa},\mathtt{b})$ & $(\eps,\mathtt{aab},\eps)$ & $\mathtt{aaaabb}$\\
$\mathtt{aababa}$ & $(\eps,\mathtt{ab},\eps)$ & $(\mathtt{a},\mathtt{aba})$ & $(\mathtt{aab},\mathtt{a})$ & $(\mathtt{a},\mathtt{b},\eps)$ & $\mathtt{aaabba}$\\
$\mathtt{aababb}$ & $(\eps,\mathtt{a},\mathtt{b})$ & $(\mathtt{aab},\mathtt{b})$ & $(\mathtt{a},\mathtt{bab})$ & $(\eps,\mathtt{ab},\eps)$ & $\mathtt{aabbab}$\\
$\mathtt{abaaab}$ & $(\mathtt{ab},\mathtt{a},\eps)$ & $(\mathtt{a},\mathtt{ab})$ & $(\mathtt{aa},\mathtt{b})$ & $(\eps,\mathtt{baa},\eps)$ & $\mathtt{aabaab}$\\
$\mathtt{ababba}$ & $(\mathtt{ab},\mathtt{b},\eps)$ & $(\mathtt{a},\mathtt{ba})$ & $(\mathtt{ab},\mathtt{a})$ & $(\eps,\mathtt{bab},\eps)$ & $\mathtt{abbaba}$\\
$\mathtt{aabaaba}$ & $(\eps,\mathtt{a},\mathtt{ba})$ & $(\mathtt{aab},\mathtt{a})$ & $(\mathtt{a},\mathtt{baa})$ & $(\eps,\mathtt{aab},\eps)$ & $\mathtt{aaabbaa}$\\
$\mathtt{abbabba}$ & $(\eps,\mathtt{b},\mathtt{bba})$ & $(\mathtt{a},\mathtt{ba})$ & $(\mathtt{ab},\mathtt{a})$ & $(\eps,\mathtt{bbab},\eps)$ & $\mathtt{abbbaba}$\\
\hline
\end{tabular}
\end{table}

\begin{lemma}[Run patterns]
\label{lem:app-singleton-runs}
Let $W\in\{\mathtt a,\mathtt b\}^+$ have length $n\ge7$ and none of the shapes of
Theorem~\ref{thm:app-singleton}\textup{(iii)}.  Then $u\sqsubseteq W$ for a word $u$
of Table~\ref{tab:singleton-certificates}.
\end{lemma}

\begin{proof}
Split $W$ into maximal runs of one letter; a run is \emph{long} if it has length at
least $2$.  The shapes are exactly the run-length sequences $(n)$, $(n-1,1)$,
$(1,n-1)$, $(1,n-2,1)$ and $(1,\ldots,1)$, so $W$ has a long run.  Since
$\sqsubseteq$ is invariant under exchange of letters and reversal, we may apply
these to $W$ in each case.

\emph{Two adjacent long runs.}  After exchange and reversal, a run $\mathtt a^p$ is
followed by a run $\mathtt b^q$ with $p,q\ge2$.  If $p\ge3$ or $q\ge3$, then
$\mathtt{aaabb}\sqsubseteq W$.  If $p=q=2$, then since $n\ge7$ the factor
$\mathtt{aabb}$ has a neighbour, $\mathtt b$ on the left or $\mathtt a$ on the right,
and $\mathtt{baabb}$ or $\mathtt{aabba}$ occurs; both are images of $\mathtt{aabba}$.

\emph{At least two long runs, no two adjacent.}  Take two long runs with the
fewest runs between them, the left one of $\mathtt a$; the $k\ge1$ runs between them
have length $1$.  If $k\ge3$,
then $\mathtt{aababa}\sqsubseteq W$; if $k=2$, then $\mathtt{aababb}\sqsubseteq W$.
If $k=1$, the two long runs form a factor
$\mathtt a^p\mathtt b\mathtt a^q$ with $p,q\ge2$.  If $p\ge3$ or $q\ge3$, then
$\mathtt{aaabaa}\sqsubseteq W$.  If $p=q=2$, at least two letters of $W$ lie
outside $\mathtt{aabaa}$; the letter next to it on either side is $\mathtt b$, and a
single one, since two long runs are not adjacent.  If $\mathtt{aabaa}$ has such a
neighbour on both sides, $W$ contains $\mathtt{baabaab}$, an image of
$\mathtt{abbabba}$.  Otherwise, say on the right, the next letter is $\mathtt a$ and
$W$ contains $\mathtt{aabaaba}$.

\emph{Exactly one long run}, say $\mathtt a^m$.  All other runs have length $1$, and since
$W$ has none of the shapes, two letters of $W$ lie on one side of $\mathtt a^m$; after
reversal, $\mathtt a^m\mathtt{ba}$ occurs.  If $m\ge4$, then
$\mathtt{aaaaba}\sqsubseteq W$.  If $m=3$, either $\mathtt{aaabab}$ occurs, or $W$
ends with $\mathtt{aaaba}$ and at least two letters precede $\mathtt a^3$, so
$\mathtt{abaaab}$ occurs.  If $m=2$, either $\mathtt{aababa}$ occurs, or $W$ ends
with $\mathtt{aabab}$ preceded by at least two letters, so $W$ contains
$\mathtt{baabab}$, or $W$ ends with $\mathtt{aaba}$ preceded by at least three
letters, so $W$ contains $\mathtt{babaab}$; both are images of
$\mathtt{ababba}$.
\end{proof}

The converse direction follows from a more general sufficient condition,
which does not depend on the size of the alphabet.

\begin{lemma}[Periodic words with marked ends]
\label{lem:app-singleton-periodic}
Let $z=z_1\cdots z_p$ be a word of $p\ge1$ pairwise distinct letters, let $w$ be a
nonempty factor of $zz\cdots$, and let $\alpha,\beta\in\Sigma\cup\{\eps\}$ consist of
letters not occurring in $z$.  Then $W:=\alpha w\beta$ passes for every arity.
\end{lemma}

\begin{proof}
Call the positions of $w$ in $W$ \emph{inner}.  An inner position carries $z_k$
with $k$ determined by the position modulo $p$, and conversely the letter $z_k$
determines the position modulo $p$.  Let $d\ge2$ and
$E[\tuple x]=E[\tuple y]=F[\tuple x]=W$ with $E=(v_0,\ldots,v_d)$ and
$F=(w_0,\ldots,w_d)$ admissible; as $E$ is common,
$\sum_i|x_i|=\sum_i|y_i|$.

\emph{End letters.}  A letter of $\alpha$ or $\beta$ occurs in $W$ only at the
first or the last position.  Since $d\ge2$ and the separators are nonempty, no
component contains both, and in $E[\tuple x]$ and $E[\tuple y]$ the first
position lies in $x_1$, respectively $y_1$, iff $v_0=\eps$; likewise at the end.
So $x_i$ and $y_i$ agree on whether they begin with $\alpha\ne\eps$ and whether
they end with $\beta\ne\eps$.

\emph{Phases.}  For $1\le j<d$ the separator $v_j$ is nonempty, does not
contain the first position, and its first letter is not at the last position, so
it starts at inner positions of $E[\tuple x]$ and $E[\tuple y]$ with the same
letter.  Their difference $\sum_{i\le j}(|y_i|-|x_i|)$ is therefore divisible by
$p$, and with the equal sums $|x_i|\equiv|y_i|\pmod p$ for all $i$.  Hence $x_i$
and $y_i$ start at positions congruent modulo $p$ and have congruent lengths; with
the end letters, they have the same first letter and the same last letter.

\emph{Conclusion.}  $W$ is the only word of length $|W|$ in which the letters of
$\alpha$ and $\beta$ occur exactly at the positions they occupy in $W$, every
other pair of adjacent letters is a pair $z_kz_{k+1}$ (indices modulo $p$), and
the letter at position $|\alpha|+1$ is that of $W$.  The word $F[\tuple y]$ has
length $|W|$.  By the first step its letters of $\alpha$ and $\beta$ sit where
they sit in $F[\tuple x]=W$.  Its other adjacent pairs lie inside some $y_i$, a
factor of $W$, inside a fixed word of $F$, or across a boundary between $F$ and
some $y_i$, where they coincide with those of $F[\tuple x]$ by the second step.
The letter at position $|\alpha|+1$ lies in $w_0$, or it is the first letter of
$y_1$ (equal to that of $x_1$), or $w_0=\eps$ and $y_1$, a prefix of $W$ in
$E[\tuple y]$, begins with $\alpha$; in each case it is that of $W$.  So
$F[\tuple y]=W$.  Arity $1$ is trivial.
\end{proof}

\begin{proof}[Proof of Theorem~\ref{thm:app-singleton}]
(ii)$\Rightarrow$(i) is trivial.

(i)$\Rightarrow$(iii).  If $W$ has none of the shapes, then
Lemma~\ref{lem:app-singleton-runs} gives a word $u$ of
Table~\ref{tab:singleton-certificates} with $u\sqsubseteq W$, and the violation for
$u$ gives one for $W$.

(iii)$\Rightarrow$(ii).  Each shape is $\alpha w\beta$ as in
Lemma~\ref{lem:app-singleton-periodic}, with $z=x$ for the first four and $z=xy$
for the alternating words.
\end{proof}

\section{Comparison with Multidimensional Substitutability}
\label{sec:comparison}

Yoshinaka's multidimensional substitutability~\cite{Yoshinaka2009,Yoshinaka2011}
is the direct positive-data baseline.  This appendix shows that trivial typing
recovers it and window typing recovers the CFG $(k,\ell)$ hierarchy
of~\cite{Yoshinaka2008}, while a fixed finite typing admits non-context-free
targets outside the baseline.
Since every $(f,s)$ condition includes arity one,
$\Cmcf{f}{s}\cap\mathrm{CFL}\subseteq\Cmcf{1}{s}$.


\begin{definition}[Multidimensional substitutability]
\label{def:yoshinaka-pd-substitutability}
A language $L$ is \emph{$f$D-substitutable} if for every $1\le d\le f$, all
$E,F\in\mathcal E_d$ and all $\tuple x,\tuple y\in(\Sigma^+)^d$,
\(E[\tuple x],E[\tuple y],F[\tuple x]\in L\Rightarrow F[\tuple y]\in L\)
(the \emph{substitution implication}); $\YSub{f}$ is the class of these languages.
This is Yoshinaka's definition, with nonempty components and nonempty internal
separators \cite[Section~3.1 and footnote~3]{Yoshinaka2011}.
\end{definition}

The substitution implication is $(f,s_{\mathrm{triv}})$-tuple substitutability,
so $\Cmcf{f}{s_{\mathrm{triv}}}=\YSub{f}\cap f\text{-}\mathrm{MCFL}\subseteq\Cmcf{f}{s}$
for every $s$.  Yoshinaka \cite[Definition~2]{Yoshinaka2008} calls $L$
\emph{$k,l$-substitutable} if for all $x_1,y_1,z_1,x_2,y_2,z_2\in\Sigma^*$,
$v\in\Sigma^k$ and $u\in\Sigma^\ell$ with $vy_1u,vy_2u\ne\eps$,
$x_1vy_1uz_1,x_1vy_2uz_1,x_2vy_1uz_2\in L$ implies $x_2vy_2uz_2\in L$.  Here
$y_1,y_2$ may be empty; only the factors $vy_iu$ must be nonempty.  Exchanging
$y_1$ and $y_2$ gives the converse implication, so the condition says that two
such factors with a common context have equal distributions.  Two nonempty words
have equal $s_{k,\ell}$-type iff they are identical or both have length at least
$\max\{1,k+\ell\}$ with the same first $k$ and last $\ell$ letters, i.e.\ iff they
are of the forms $vy_1u,vy_2u$ above, with $y_1,y_2$ possibly empty, or equal.
Words shorter than $k+\ell$ are therefore compared only with themselves, in line
with Yoshinaka's remark that the condition says nothing about them.  Hence
$(1,s_{k,\ell})$-tuple substitutability is exactly Yoshinaka's $(k,\ell)$
condition; the componentwise version at higher fan-out is a definition of this
paper.

A language $L$ is \emph{strictly 2-local} if there is a set $\Pi$ of pairs over
$\Sigma\cup\{\vdash,\dashv\}$ such that a word $w$ lies in $L$ iff every pair of
adjacent symbols of $\vdash w\dashv$ lies in $\Pi$.

\begin{lemma}[Strictly 2-local languages are boundary typed]
\label{lem:sl2-boundary-control}
Every strictly 2-local $L$ is $(f,s_{1,1})$-tuple-substitutable for every $f$; in
particular $L\in\Cmcf{f}{s_{1,1}}$, and $L\in\Cmcf{f}{s}$ whenever $s_{1,1}\preceq s$.
\end{lemma}

\begin{proof}
Let $\tuple x,\tuple y\in(\Sigma^+)^d$ have equal $s_{1,1}$-types, so
corresponding components are identical letters or share first and last letters,
and let $E_0[\tuple x],E_0[\tuple y]\in L$.  Suppose $E[\tuple x]\in L$.  Since the
internal separators of $E$ are nonempty, every adjacent pair of
$\vdash E[\tuple y]\dashv$ either lies inside some $y_i$, and then occurs in
$E_0[\tuple y]\in L$, or joins a first or last letter of some $y_i$ to a symbol of
$\vdash E\dashv$, or lies in $\vdash E\dashv$, and then occurs in
$\vdash E[\tuple x]\dashv$.  So $E[\tuple y]\in L$; symmetry and $s_{1,1}\preceq s$
finish the proof (regular languages are MCFLs).
\end{proof}

\begin{proposition}[Cross-serial dependencies and feature licensing]
\label{prop:feature-cross-typed}
Let $\Sigma_{\mathrm{feat}}=\{\mathtt p_0,\mathtt p_1,\allowbreak\mathtt q_0,\mathtt q_1,\allowbreak
\mathtt a,\mathtt b,\mathtt c,\mathtt d\}$, let $g:\Sigma_{\mathrm{feat}}^*\to\mathbb Z_2^2$
map $\mathtt p_1\mapsto(1,0)$, $\mathtt q_1\mapsto(0,1)$ and all other letters to $0$, and
let $s:=(s_{1,1},g)|_{\Sigma_{\mathrm{feat}}^+}$ (corestricted to its image).  For
$R\subseteq\{0,1\}^2$ put
\[
 L_R:=\{\mathtt a^m\mathtt p_i\mathtt b^n\mathtt c^m\mathtt q_j\mathtt d^n:
 m,n\ge1,\ (i,j)\in R\}.
\]
Then for every $R$ and every $f\ge2$, $L_R\in\mathsf{MCFL}(2,2)\cap\Cmcf{f}{s}$, and
$L_R$ is not context-free if $R\ne\emptyset$.  For
$\mathcal R_{\mathrm{feat}}=\{(0,0),(0,1),(1,0)\}$, the language
$L_{\mathcal R_{\mathrm{feat}}}$ is neither in $\YSub{f}$ nor
$(1,s_{1,1})$-tuple-substitutable.  The typing is fixed once for the whole
family.
\end{proposition}

\begin{proof}
\emph{Grammar.}  With fan-out-two $A,B$ and rules
$A\to(\mathtt a,\mathtt c)\mid(\mathtt a\xi_1^1,\mathtt c\xi_1^2)(A)$,
$B\to(\mathtt b,\mathtt d)\mid(\mathtt b\xi_1^1,\mathtt d\xi_1^2)(B)$ and
$S\to(\xi_1^1\mathtt p_i\xi_2^1\xi_1^2\mathtt q_j\xi_2^2)(A,B)$ for
$(i,j)\in R$, we get $L_A=\{(\mathtt a^m,\mathtt c^m)\}$,
$L_B=\{(\mathtt b^n,\mathtt d^n)\}$ and $L(G)=L_R$; the grammar is good,
of fan-out and rank two.  For $(i,j)\in R$, intersecting with
$\mathtt a^+\mathtt p_i\mathtt b^+\mathtt c^+\mathtt q_j\mathtt d^+$ and erasing the
markers gives the non-context-free $\{\mathtt a^m\mathtt b^n\mathtt c^m\mathtt d^n\}$.

\emph{Typed substitutability.}  Let
$Q=\mathtt a^+\{\mathtt p_0,\mathtt p_1\}\mathtt b^+\mathtt c^+\{\mathtt q_0,\mathtt q_1\}
\mathtt d^+$ (strictly 2-local), $K_R:=\{w\in\Sigma_{\mathrm{feat}}^+:g(w)\in R\}$, and let
$\varphi$ map $\mathtt a,\mathtt c,\mathtt b,\mathtt d$ to
$(1,0),(-1,0),(0,1),(0,-1)\in\mathbb Z^2$ and the markers to $0$.  On $Q$,
$g(w)=(i,j)$ for the markers $\mathtt p_i,\mathtt q_j$ of $w$, so
$L_R=Q\cap K_R\cap\varphi^{-1}(0)$.  The three languages are
$(f,s)$-tuple-substitutable by Lemma~\ref{lem:sl2-boundary-control}, by item~(c)
after Definition~\ref{def:class-cmcf} ($K_R$ is a union of $s$-classes), and by
Proposition~\ref{prop:abelian-coset-filter}(ii);
Proposition~\ref{prop:abelian-coset-filter}(i) gives the intersection.

\emph{Failure without $g$.}  For $R=\mathcal R_{\mathrm{feat}}$, $x=\mathtt{ap_0b}$,
$y=\mathtt{ap_1b}$, $E=\blank_1\mathtt{cq_0d}$ and $F=\blank_1\mathtt{cq_1d}$ we have
$E[x],E[y],F[x]\in L_R$ but $F[y]\notin L_R$, since $(1,1)\notin R$.  Here
$s_{1,1}(x)=s_{1,1}(y)$ but $g(x)\ne g(y)$.  The same four words violate the
substitution implication of Definition~\ref{def:yoshinaka-pd-substitutability} at
$d=1$, so $L_R\notin\YSub{f}$ for every $f\ge1$.
\end{proof}


